\documentclass[11pt]{amsart}

% Standard packages only (all in every TeX Live / MiKTeX installation used by arXiv).
\usepackage[margin=1.05in]{geometry}
\usepackage{amssymb}
\usepackage{alltt}
\usepackage{booktabs}
\usepackage[hidelinks]{hyperref}

% For arXiv: run pdflatex, bibtex, pdflatex, pdflatex and upload main.bbl together with main.tex
% (a hand-made main.bbl consistent with refs.bib is provided in this folder).

\newtheorem{mainthm}{Theorem}
\newtheorem{lemma}{Lemma}[section]
\newtheorem{proposition}[lemma]{Proposition}
\newtheorem{corollary}[lemma]{Corollary}
\theoremstyle{definition}
\newtheorem{definition}[lemma]{Definition}
\theoremstyle{remark}
\newtheorem{remark}[lemma]{Remark}

\numberwithin{equation}{section}

\newcommand{\R}{\mathbb{R}}
\newcommand{\C}{\mathbb{C}}
\newcommand{\N}{\mathbb{N}}
\newcommand{\Tr}{\operatorname{Tr}}
\newcommand{\diag}{\operatorname{diag}}
\newcommand{\conv}{\operatorname{conv}}
\newcommand{\Reach}{\operatorname{Reach}}
\newcommand{\row}{\operatorname{row}}
\newcommand{\run}{\operatorname{run}}
\newcommand{\tail}{\operatorname{tail}}
\newcommand{\ind}[1]{\mathbf{1}[#1]}
\newcommand{\dom}{\trianglelefteq}
\newcommand{\sT}{\mathsf{T}}
\newcommand{\sH}{\mathsf{H}}
\newcommand{\sG}{\mathsf{G}}
\newcommand{\cB}{\mathcal{B}}
\newcommand{\cQ}{\mathcal{Q}}
\newcommand{\cR}{\mathcal{R}}
\newcommand{\ms}{\mathfrak{m}}
\newcommand{\aup}{a^{+}}
\newcommand{\alo}{a^{-}}
\newcommand{\lean}[1]{\texttt{#1}}

\setlength{\emergencystretch}{3em}
\begin{document}

\title[Symmetrically thermalizing unitaries]{Symmetrically thermalizing unitaries exist in every dimension}

\author{Ansh Mishra}
\author{Aryan Senthilkumar}

\subjclass[2020]{81P40; 15A42; 15B51; 68V20}
\keywords{quantum thermodynamics, correlations, quantum marginal problem, majorization,
Schur--Horn theorem, formal verification, Lean}

\begin{abstract}
Two identical quantum systems with Hamiltonian $H$, both in the thermal state $\tau_\beta$, can be
correlated by a global unitary. The energy needed to create a given amount of mutual information is
bounded from below, and the bound is attained when there is a \emph{symmetrically thermalizing
unitary} (STU): a unitary $U$ on $\C^d\otimes\C^d$ such that both marginals of
$U(\tau_\beta\otimes\tau_\beta)U^\dagger$ are thermal at one common higher temperature,
$\tau_{\beta'}$ with $\beta'\le\beta$. Bakhshinezhad, Clivaz, Vitagliano, Erker, Rezakhani, Huber and
Friis (2019) proved that STUs exist in local dimensions $d\le 4$ and conjectured that they exist in
every dimension. We prove this conjecture: for every $d$, every Hamiltonian $H$, every $\beta>0$ and
every $\beta'\in[0,\beta]$ a symmetrically thermalizing unitary exists. In an eigenbasis of $H$ it can
be chosen real orthogonal, and it is given by an explicit construction. The proof is elementary and
does not use Kirwan's convexity theorem. It reduces the problem to majorization conditions on the
diagonals of a $d\times d$ array, writes the target as a convex combination of clipped thermal
vectors, and builds, for every interval of hook lengths and every box size, an explicit symmetric
``interval uniformisation'' matrix from renewal squares and at most one correcting dipole; the
dipole amplitude is fixed by identities for alternating harmonic sums. The complete proof, for
every Hermitian $H$ (via the case $H=\diag(E)$ with sorted energies and a local unitary), has been
formalised in Lean~4 with Mathlib and uses only Lean's standard axioms. We also show that the
stronger statement ``every spectrum majorized by the initial one is symmetrically reachable'' is
false in every dimension $d\ge 4$.
\end{abstract}

\maketitle

\tableofcontents

%==================================================================================================
\section{Introduction}\label{sec:intro}
%==================================================================================================

\subsection{Correlations at minimal energy cost}
Correlations between physical systems are a resource, and creating them costs energy
\cite{HuberEtAl2015,BruschiEtAl2015,VitaglianoEtAl2019}. A basic instance of this trade-off is the
following \cite{BakhshinezhadEtAl2019}. Two identical quantum systems $A$ and $B$, each with
Hamiltonian $H$ on $\C^d$, are initially uncorrelated and in equilibrium with a bath at inverse
temperature $\beta>0$, so that the global state is $\tau_\beta\otimes\tau_\beta$ with
\[
  \tau_\beta=\tau_\beta(H):=\frac{e^{-\beta H}}{\Tr e^{-\beta H}} .
\]
This initial state is completely passive \cite{PuszWoronowicz1978}: no energy can be extracted from it
by unitaries, even from many copies.
With full control of the closed composite system one may apply any global unitary $U$, producing
$\rho'=U(\tau_\beta\otimes\tau_\beta)U^\dagger$. The invested energy
$\Delta E=\Tr[(H\otimes 1+1\otimes H)(\rho'-\tau_\beta\otimes\tau_\beta)]$ depends only on the
marginals $\rho'_A=\Tr_B\rho'$ and $\rho'_B=\Tr_A\rho'$. Measuring correlations by the mutual
information $I(A{:}B)=S(\rho_A)+S(\rho_B)-S(\rho_{AB})$, and using that the global von Neumann
entropy is invariant under unitaries, the created correlations are $\Delta I=\Delta S_A+\Delta S_B$.
Since thermal states maximise the entropy at fixed energy, and the entropy of thermal states is a
concave function of their energy, an invested energy $\Delta E=2\Tr[H(\tau_{\beta'}-\tau_\beta)]$ with
$0\le\beta'\le\beta$ can create at most $\Delta I\le2\bigl[S(\tau_{\beta'})-S(\tau_\beta)\bigr]$, with
equality if both final marginals equal $\tau_{\beta'}$ \cite[Sec.~II and App.~A.I]{BakhshinezhadEtAl2019}.
The resulting lower bound on the energy cost of correlations is therefore attained whenever the
following unitaries exist.

\begin{definition}\label{def:stu}
Let $H$ be a Hamiltonian on $\C^d$ and $0\le\beta'\le\beta$. A unitary $U$ on $\C^d\otimes\C^d$ is
\emph{symmetrically thermalizing} (an STU) for $(H,\beta,\beta')$ if
\[
  \Tr_B\bigl[U(\tau_\beta\otimes\tau_\beta)U^\dagger\bigr]
  =\Tr_A\bigl[U(\tau_\beta\otimes\tau_\beta)U^\dagger\bigr]=\tau_{\beta'} .
\]
\end{definition}

Bakhshinezhad, Clivaz, Vitagliano, Erker, Rezakhani, Huber and Friis asked whether STUs exist for
every Hamiltonian, every $\beta$ and every $\beta'\le\beta$ \cite[Question~2]{BakhshinezhadEtAl2019},
and conjectured that they do \cite[Sec.~IV]{BakhshinezhadEtAl2019}. The symmetry of the problem is
essential: when the two systems have different Hamiltonians the analogous statement fails
\cite[Sec.~II.1.1]{BakhshinezhadEtAl2019}.

\subsection{Previous status}
STUs were known to exist when the spectrum of $H$ is equally spaced \cite{HuberEtAl2015} (in
particular for qubits, $d=2$), and for every Hamiltonian in local dimensions $d=3$ and $d=4$
\cite{BakhshinezhadEtAl2019}. For $d\ge5$, Ref.~\cite{BakhshinezhadEtAl2019} gave sufficient
conditions and a framework of locally classical subspaces; the general case was still open in the
PhD thesis \cite{Clivaz2020}. See also \cite{JevticJenningsRudolph2012,McKayEtAl2018} for related
questions about correlations in unitary orbits.

\subsection{Main result}
We prove the conjecture.

\medskip
\noindent\textbf{Theorem} (informal version of Theorems~\ref{thm:main} and~\ref{thm:diag}).
\emph{For every local dimension $d$, every Hamiltonian $H$ on $\C^d$, every $\beta>0$ and every
$\beta'\in[0,\beta]$ there is a symmetrically thermalizing unitary. In an eigenbasis of $H$ it can be
chosen to be a real orthogonal matrix, and it is given by an explicit construction.}
\medskip

Three features of the proof deserve mention.
\begin{itemize}
\item \emph{Elementary and explicit.} The only tools are majorization, the existence direction of
the real Schur--Horn theorem (in its constructive form \cite{ChanLi1983}), and elementary estimates
for alternating series. Every object in the construction is explicit.
\item \emph{Kirwan-free.} For a fixed global spectrum, the set of pairs of marginal spectra in the
unitary orbit is a convex polytope by Kirwan's convexity theorem \cite{Kirwan1984}; see
\cite{Klyachko2004,DaftuarHayden2005,ChristandlMitchison2006} for this quantum marginal problem.
Convexity reduces the conjecture to the reachability of finitely many explicit points, which gives an
alternative route that we recall in Section~\ref{sec:kirwan}. Our proof does not use it: we work
inside one fixed combinatorial structure in which convexity is automatic.
\item \emph{Machine-checked.} The complete proof of Theorem~\ref{thm:main} has been formalised in
Lean~4 \cite{deMouraUllrich2021} with Mathlib \cite{mathlib2020}: Theorem~\ref{thm:diag}, the statement
for $H=\diag(E)$ with sorted energies, and the reduction of an arbitrary Hermitian $H$ to it, with
$\tau_\beta(H)=e^{-\beta H}/\Tr e^{-\beta H}$ defined by the matrix exponential. This includes all
analytic estimates. The formalisation uses only Lean's three standard axioms and contains no
unproved statements (Section~\ref{sec:lean}).
\end{itemize}

\subsection{Strategy}
Let $p_0\ge p_1\ge\dots\ge p_{d-1}>0$ be the eigenvalues of $\tau_\beta$ and $q$ those of
$\tau_{\beta'}$, in the same eigenbasis of $H$.
\begin{enumerate}
\item \emph{Classical reduction} (Section~\ref{sec:reduction}). Arrange the product eigenvectors
$|i\rangle|j\rangle$ of $\tau_\beta\otimes\tau_\beta$ in a $d\times d$ array. Each of its $2d-1$
diagonals meets every row and every column at most once, so a Schur--Horn rotation inside a diagonal
keeps both marginals diagonal. Hence it suffices to find a real symmetric $d\times d$ matrix $F$
whose diagonals are majorized by the corresponding diagonals of $p\otimes p$ and whose row sums are
$q$. The set $\Reach(p)$ of such row-sum vectors is convex.
\item \emph{Clipped points} (Lemma~\ref{lem:T}). The vector $q$ is a convex combination of the
clipped thermal vectors $w_k=\min(p,p_k)/Z_k$, $0\le k\le d-1$. This holds for every concave
nondecreasing reweighting of $p$, in particular for $q\propto p^{\beta'/\beta}$.
\item \emph{Young diagrams} (Sections~\ref{sec:reduction} and~\ref{sec:unif}). With
$B=\min(p,p_k)$, the difference $p\otimes p-B\otimes B$ is a nonnegative combination of symmetric Young
diagrams (squares and crosses). All pieces are nonincreasing along every diagonal, so majorizations
can be superposed. Each piece is replaced by an \emph{interval uniformisation}: a symmetric matrix
with entries in $[0,1]$, the same diagonal sums as the piece, and constant row sums. The pieces are
distributed over square boxes according to an explicit capacity (Hall-type) identity. This gives
$w_k\in\Reach(p)$.
\item \emph{Interval uniformisations} (Sections~\ref{sec:unif}--\ref{sec:cases}). We construct them
for every box size $n$ and every interval $[\ell,h]$ of hook lengths: renewal squares, their
complements, and a Toeplitz-plus-Hankel matrix (the \emph{T/H/G matrix}) corrected by at most one
\emph{dipole}. Diagonal sums and row sums hold by exact identities. Nonnegativity and the bound $1$
reduce to one-dimensional inequalities, which follow from six facts F1--F6 about the numbers
\[
  a(K)=4K^2(-1)^K\Bigl(\ln 2-\sum_{i=1}^{K}\frac{(-1)^{i-1}}{i}\Bigr)+1-2K ,
\]
proved with a two-sided bound on the tails of the alternating harmonic series.
\end{enumerate}
Finally, the ``strong form'' of the problem, in which every spectrum majorized by $p$ is
symmetrically reachable, holds for $d=3$ \cite[Thm.~1]{BakhshinezhadEtAl2019} but fails for every
$d\ge4$ (Proposition~\ref{prop:counterexample}). So the thermal form of the target must be used, as it
is in step~(2).

\subsection{Organisation}
Section~\ref{sec:statement} states the results precisely, including the exact Lean statement, and
reduces a general Hamiltonian to the diagonal sorted case. Section~\ref{sec:reduction} contains the
classical reduction, Lemma~T and the decomposition of $p\otimes p$. Section~\ref{sec:unif} defines
interval uniformisations, proves that they suffice, and treats renewal squares and complements.
Sections~\ref{sec:thg}--\ref{sec:cases} construct the remaining interval uniformisations.
Section~\ref{sec:lean} describes the formal verification and Section~\ref{sec:remarks} contains
further remarks: the counterexample to the strong form, the Kirwan route, and an outlook.
Appendix~\ref{app:poly} lists the polynomial identities behind F1--F4, and Appendix~\ref{app:checks}
lists the independent computer checks of the formulas in this paper.

%==================================================================================================
\section{Statement of the results}\label{sec:statement}
%==================================================================================================

We write $[d]=\{0,1,\dots,d-1\}$ and use the standard basis $|i\rangle$, $i\in[d]$, of $\C^d$.
Operators on $\C^d\otimes\C^d$ are matrices indexed by $[d]\times[d]$, and the partial traces are
\[
  (\Tr_B X)_{a,a'}=\sum_{b\in[d]}X_{(a,b),(a',b)},\qquad
  (\Tr_A X)_{b,b'}=\sum_{a\in[d]}X_{(a,b),(a,b')} .
\]
For $E\in\R^d$ and $\beta\in\R$ the \emph{Gibbs vector} is
$\gamma_\beta(E)_i=e^{-\beta E_i}/\sum_{j}e^{-\beta E_j}$, so that
$\tau_\beta(\diag E)=\diag\gamma_\beta(E)$.

\begin{mainthm}[STUs exist in every dimension]\label{thm:main}
Let $d\ge1$, let $H$ be a Hermitian operator on $\C^d$, and let $\beta>0$ and $0\le\beta'\le\beta$.
There is a unitary $U$ on $\C^d\otimes\C^d$ such that
\[
  \Tr_B\bigl[U(\tau_\beta(H)\otimes\tau_\beta(H))U^\dagger\bigr]
  =\Tr_A\bigl[U(\tau_\beta(H)\otimes\tau_\beta(H))U^\dagger\bigr]=\tau_{\beta'}(H) .
\]
If $H=V\diag(E)V^\dagger$ with $V$ unitary and $E_0\le\dots\le E_{d-1}$, one can take
$U=(V\otimes V)\,O\,(V\otimes V)^\dagger$ with $O$ a real orthogonal matrix.
\end{mainthm}

\begin{mainthm}[diagonal form]\label{thm:diag}
Let $d\in\N$, let $E\in\R^d$ with $E_0\le E_1\le\dots\le E_{d-1}$, let $\beta>0$ and
$0\le\beta'\le\beta$, and put $D_\beta=\diag\gamma_\beta(E)$. There is a unitary $U$ on
$\C^d\otimes\C^d$, which can be taken to be a real orthogonal matrix, such that
\[
  \Tr_B\bigl[U(D_\beta\otimes D_\beta)U^\dagger\bigr]=\Tr_A\bigl[U(D_\beta\otimes D_\beta)U^\dagger\bigr]
  =D_{\beta'} .
\]
\end{mainthm}

The cases $\beta'=\beta$ ($U=1$) and $d\le1$ are trivial, and for $\beta=0$ necessarily $\beta'=0$ and
$U=1$ works. Degenerate energies are allowed.

\begin{remark}[the formal statement]\label{rem:leanstatement}
Except for the clause ``which can be taken to be a real orthogonal matrix'', Theorem~\ref{thm:diag}
is literally the Lean theorem \lean{STUProof.stu\_exists\_unconditional} in the file
\lean{lean/STUProof/KirwanFreeFinal.lean} of the repository \cite{STUrepo}:
\begin{alltt}\small
theorem stu_exists_unconditional :
    \(\forall\) (d : \(\mathbb{N}\)) (E : Fin d \(\to\) \(\mathbb{R}\)), Monotone E \(\to\)
    \(\forall\) \(\beta\) \(\beta\sp{\prime}\) : \(\mathbb{R}\), 0 < \(\beta\) \(\to\) 0 \(\le\) \(\beta\sp{\prime}\) \(\to\) \(\beta\sp{\prime}\) \(\le\) \(\beta\) \(\to\)
      \(\exists\) U \(\in\) unitaryGroup (Fin d \(\times\) Fin d) \(\mathbb{C}\),
        partialTraceB (U * (gibbsState E \(\beta\) \(\otimes\sb{k}\) gibbsState E \(\beta\)) * star U)
            = gibbsState E \(\beta\sp{\prime}\) \(\wedge\)
        partialTraceA (U * (gibbsState E \(\beta\) \(\otimes\sb{k}\) gibbsState E \(\beta\)) * star U)
            = gibbsState E \(\beta\sp{\prime}\)
\end{alltt}
Here \lean{Monotone E} means $E_0\le\dots\le E_{d-1}$, \lean{star U} is $U^\dagger$,
$\otimes_k$ is the Kronecker product, and the four project definitions entering the statement are
\begin{alltt}\small
def gibbs (E : Fin d \(\to\) \(\mathbb{R}\)) (\(\beta\) : \(\mathbb{R}\)) : Fin d \(\to\) \(\mathbb{R}\) :=
  fun i => Real.exp (-\(\beta\) * E i) / \(\sum\) j, Real.exp (-\(\beta\) * E j)
def gibbsState (E : Fin d \(\to\) \(\mathbb{R}\)) (\(\beta\) : \(\mathbb{R}\)) : Matrix (Fin d) (Fin d) \(\mathbb{C}\) :=
  diagonal fun i => (gibbs E \(\beta\) i : \(\mathbb{C}\))
def partialTraceB (\(\rho\) : Matrix (m \(\times\) n) (m \(\times\) n) R) : Matrix m m R :=
  Matrix.of fun a a\(\sp{\prime}\) => \(\sum\) b, \(\rho\) (a, b) (a\(\sp{\prime}\), b)
def partialTraceA (\(\rho\) : Matrix (m \(\times\) n) (m \(\times\) n) R) : Matrix n n R :=
  Matrix.of fun b b\(\sp{\prime}\) => \(\sum\) a, \(\rho\) (a, b) (a, b\(\sp{\prime}\))
\end{alltt}
The proof constructs $U$ as a real orthogonal matrix, embedded in the complex matrices.
\end{remark}

\begin{proof}[Proof of Theorem~\ref{thm:main} from Theorem~\ref{thm:diag}]
By the spectral theorem, $H=V\diag(E)V^\dagger$ with $V$ unitary and $E_0\le\dots\le E_{d-1}$ (list the
eigenvalues in nondecreasing order). Then $e^{-\beta H}=V\diag(e^{-\beta E_i})V^\dagger$, hence
$\tau_\beta(H)=VD_\beta V^\dagger$ for every $\beta$. For unitaries $V,W$ on $\C^d$ and every operator
$X$ on $\C^d\otimes\C^d$,
\begin{equation}\label{eq:localunitary}
  \Tr_B\bigl[(V\otimes W)X(V\otimes W)^\dagger\bigr]=V(\Tr_B X)V^\dagger,\qquad
  \Tr_A\bigl[(V\otimes W)X(V\otimes W)^\dagger\bigr]=W(\Tr_A X)W^\dagger .
\end{equation}
Indeed, the $(a,a')$ entry of the first left-hand side is
$\sum_{b,x,y,x',y'}V_{ax}W_{by}X_{(x,y),(x',y')}\overline{V_{a'x'}}\,\overline{W_{by'}}$, and
$\sum_b W_{by}\overline{W_{by'}}=\delta_{yy'}$ because $W^\dagger W=1$; the second identity is the same
computation with the roles of the factors exchanged. Let $O$ be the real orthogonal matrix given by
Theorem~\ref{thm:diag} and put $U=(V\otimes V)O(V\otimes V)^\dagger$. Then
$U(\tau_\beta(H)\otimes\tau_\beta(H))U^\dagger=(V\otimes V)\,O(D_\beta\otimes D_\beta)O^{\mathsf T}(V\otimes V)^\dagger$,
and by \eqref{eq:localunitary} both marginals equal $VD_{\beta'}V^\dagger=\tau_{\beta'}(H)$.
\end{proof}

The identities \eqref{eq:localunitary} are formalised (\lean{partialTraceB\_conj\_kronecker} and
\lean{partialTraceA\_conj\_kronecker} in \lean{STU.lean}). The reduction itself is also formalised: the
theorem \lean{stu\_exists\_hermitian} (\lean{GeneralH.lean}) states Theorem~\ref{thm:main} for every
Hermitian $H$ on $\C^d$, with $\tau_\beta(H)=e^{-\beta H}/\Tr e^{-\beta H}$ defined through Mathlib's matrix
exponential, and derives it from Theorem~\ref{thm:diag} via the spectral theorem with sorted eigenvalues,
the covariance of $\tau_\beta$ under unitary conjugation, and the local unitary $V\otimes V$. Hence
Theorem~\ref{thm:main} itself is machine-checked.

The rest of the paper proves Theorem~\ref{thm:diag}. The skeleton of the proof is:
\begin{itemize}
\item Proposition~\ref{prop:reduction}: Theorem~\ref{thm:diag} holds in dimension $d$ if interval
uniformisations exist for all intervals of all boxes of size $n\le d$ (Lean:
\lean{stu\_exists\_of\_intervalUnif});
\item Theorem~\ref{thm:unif}: interval uniformisations exist for every $n$ (Lean:
\lean{intervalUnifAll\_all}).
\end{itemize}

%==================================================================================================
\section{Reduction to a classical problem}\label{sec:reduction}
%==================================================================================================

\subsection{Majorization}
For $x\in\R^\ell$ let $x^\downarrow$ be its nonincreasing rearrangement. Recall that $x$ is
\emph{majorized} by $\lambda$, written $x\prec\lambda$, if $\sum_i x_i=\sum_i\lambda_i$ and
$\sum_{i<k}x^\downarrow_i\le\sum_{i<k}\lambda^\downarrow_i$ for $1\le k\le\ell$ \cite{MarshallOlkinArnold2011}.
We use the following variant, which does not sort $\lambda$.

\begin{definition}\label{def:dom}
For $x,\lambda\in\R^\ell$ we write $x\dom\lambda$ if $\sum_{i<\ell}x_i=\sum_{i<\ell}\lambda_i$ and
\[
  \sum_{i\in S}x_i\le\sum_{i<|S|}\lambda_i\qquad\text{for every }S\subseteq[\ell] .
\]
\end{definition}

\begin{lemma}\label{lem:dom}
Let $x,x',\lambda,\lambda'\in\R^\ell$.
\begin{enumerate}
\item[(a)] If $x\dom\lambda$ then $x\prec\lambda$. If $\lambda$ is nonincreasing, then
$x\dom\lambda$ if and only if $x\prec\lambda$.
\item[(b)] \textup{(Superposition.)} If $x\dom\lambda$ and $x'\dom\lambda'$, then
$x+x'\dom\lambda+\lambda'$; if $c\ge0$ then $cx\dom c\lambda$.
\item[(c)] If $\lambda$ is nonincreasing, then $\lambda\dom\lambda$.
\item[(d)] \textup{(0/1 lemma.)} If $0\le x_i\le1$ for all $i$ and $\sum_i x_i=N\in\{0,\dots,\ell\}$,
then $x\dom\mathbf 1_N$, where $\mathbf 1_N=(1,\dots,1,0,\dots,0)$ has $N$ ones.
\end{enumerate}
\end{lemma}

\begin{proof}
(a) If $S$ indexes $k$ largest entries of $x$, then
$\sum_{i<k}x^\downarrow_i=\sum_{i\in S}x_i\le\sum_{i<k}\lambda_i\le\sum_{i<k}\lambda^\downarrow_i$.
If $\lambda$ is nonincreasing and $x\prec\lambda$, then for $|S|=k$,
$\sum_{i\in S}x_i\le\sum_{i<k}x^\downarrow_i\le\sum_{i<k}\lambda^\downarrow_i=\sum_{i<k}\lambda_i$.
(b) Add, respectively scale, the defining inequalities. (c) is the case $x=\lambda$ of (a).
(d) For every $S$, $\sum_{i\in S}x_i\le\min(|S|,N)=\sum_{i<|S|}(\mathbf 1_N)_i$.
\end{proof}

In the form ``if $x^{(m)}\prec\lambda^{(m)}$ and all $\lambda^{(m)}$ are nonincreasing, then
$\sum_m x^{(m)}\prec\sum_m\lambda^{(m)}$'', part (b) is the superposition principle used throughout.
The common ordering matters: $(0,1)\prec(1,0)$ and $(0,1)\prec(0,1)$, but $(0,2)\not\prec(1,1)$.

\subsection{Configurations and their realisation}
Throughout Sections~\ref{sec:reduction} and~\ref{sec:unif}, $p=(p_0,\dots,p_{d-1})$ satisfies
\[
  p_0\ge p_1\ge\dots\ge p_{d-1}>0,\qquad \sum_{i}p_i=1 .
\]

\begin{definition}\label{def:config}
A \emph{configuration} is a real symmetric $d\times d$ matrix $F=(F(i,j))_{i,j\in[d]}$. For
$0\le s\le d-1$ its \emph{$s$-th run} is $F^{(s)}=(F(i,i+s))_{0\le i\le d-1-s}\in\R^{d-s}$. The runs
of $p\otimes p$ are $\lambda^{(s)}=(p_ip_{i+s})_{0\le i\le d-1-s}$; they are nonincreasing. $F$ is
\emph{admissible} for $p$ if $F^{(s)}\dom\lambda^{(s)}$ for every $s$. We put
\[
  \Reach(p)=\Bigl\{\bigl(\textstyle\sum_{j}F(i,j)\bigr)_{i\in[d]} : F\text{ admissible for }p\Bigr\}\subseteq\R^d .
\]
\end{definition}

The key classical input is the existence direction of the Schur--Horn theorem, due to Horn
\cite{Schur1923,Horn1954}; a constructive proof by Givens rotations is in \cite{ChanLi1983}.

\begin{lemma}[Horn]\label{lem:horn}
If $x,\lambda\in\R^N$ and $x\prec\lambda$, there is a real orthogonal $N\times N$ matrix $G$ such
that the diagonal of $G\diag(\lambda)G^{\mathsf T}$ is $x$.
\end{lemma}

This lemma is formalised in the project (\lean{schurHorn} in \lean{SchurHorn.lean}), with a
Chan--Li type induction and the ``hockey-stick'' form of majorization,
$\sum_i\max(x_i-t,0)\le\sum_i\max(\lambda_i-t,0)$ for all $t\in\R$, which needs no sorting.

\begin{proposition}[realisation]\label{prop:realisation}
Let $F$ be admissible for $p$ and $q_i=\sum_jF(i,j)$. There is a real orthogonal matrix $O$ on
$\R^{[d]\times[d]}$ such that
\[
  \Tr_B\bigl[O\diag(p\otimes p)O^{\mathsf T}\bigr]=\Tr_A\bigl[O\diag(p\otimes p)O^{\mathsf T}\bigr]=\diag(q) ,
\]
where $\diag(p\otimes p)$ has entries $p_ip_j$ at $(i,j)$. The matrix $O$ is block diagonal with respect
to the partition of $[d]\times[d]$ into the $2d-1$ diagonals $D_t=\{(i,j):j-i=t\}$, $|t|\le d-1$ (the
loop $t=0$, the upper runs $t>0$ and the lower runs $t<0$; this is the ``tri-run structure'' of the
Lean development).
\end{proposition}

\begin{proof}
Each diagonal $D_t$ meets every row and every column of the array at most once. For $t=s\ge0$ the
vector $(F(i,j))_{(i,j)\in D_s}$ is the run $F^{(s)}$, and $(p_ip_j)_{(i,j)\in D_s}=\lambda^{(s)}$; for
$t=-s<0$ the two vectors are the same, by the symmetry of $F$ and of $p\otimes p$. By admissibility
and Lemma~\ref{lem:dom}(a), $F|_{D_t}\prec(p\otimes p)|_{D_t}$, so Lemma~\ref{lem:horn} gives a real
orthogonal $G_t$ on $\R^{D_t}$ with $\diag(G_t\diag((p\otimes p)|_{D_t})G_t^{\mathsf T})=F|_{D_t}$. Let
$O=\bigoplus_tG_t$ and $\rho=O\diag(p\otimes p)O^{\mathsf T}=\bigoplus_tG_t\diag((p\otimes p)|_{D_t})G_t^{\mathsf T}$.
Then $\rho_{(i,j),(i,j)}=F(i,j)$, and $\rho_{(i,j),(i',j')}=0$ unless $(i,j)$ and $(i',j')$ lie in the same
diagonal. In $(\Tr_B\rho)_{a,a'}=\sum_b\rho_{(a,b),(a',b)}$ with $a\ne a'$, the cells $(a,b)$ and $(a',b)$
lie in the same column, hence in different diagonals, so every term vanishes. Therefore
$\Tr_B\rho=\diag(\sum_bF(a,b))_a=\diag(q)$. In the same way
$\Tr_A\rho=\diag(\sum_aF(a,b))_b$, which is $\diag(q)$ because $F$ is symmetric.
\end{proof}

\begin{lemma}[convexity]\label{lem:convex}
$\Reach(p)$ is convex.
\end{lemma}

\begin{proof}
Let $F_k$ be admissible with row sums $q^{(k)}$, and let $\mu_k\ge0$ with $\sum_k\mu_k=1$. Then
$F=\sum_k\mu_kF_k$ is symmetric with row sums $\sum_k\mu_kq^{(k)}$, and by Lemma~\ref{lem:dom}(b),
$F^{(s)}=\sum_k\mu_kF_k^{(s)}\dom\sum_k\mu_k\lambda^{(s)}=\lambda^{(s)}$.
\end{proof}

In Lean, admissible configurations and $\Reach$ are \lean{ReachTri}, Lemma~\ref{lem:dom} is
\lean{SMaj} with \lean{SMaj.add}, \lean{SMaj.smul}, \lean{SMaj.majOn}, \lean{smaj\_self},
\lean{smaj\_zeroOne}, Proposition~\ref{prop:realisation} is \lean{ReachTri.orthogonal} (via
\lean{tri\_reachable} and \lean{lc\_reachable}), and Lemma~\ref{lem:convex} is \lean{ReachTri.convex};
all are in \lean{KirwanFreeCore.lean}.

\subsection{Lemma T: thermal targets are mixtures of clipped points}

\begin{lemma}[Lemma T]\label{lem:T}
Let $P_0\ge P_1\ge\dots\ge P_{d-1}>0$ and let $g:[0,\infty)\to\R$ be concave and nondecreasing with
$g(0)\ge0$. There are $\nu_0,\dots,\nu_{d-1}\ge0$ such that
\[
  g(P_i)=\sum_{k=0}^{d-1}\nu_k\min(P_i,P_k)\qquad\text{for all }i\in[d] .
\]
\end{lemma}

\begin{proof}
Let $N_0=0$. For $1\le k\le d-1$ let $N_k=\frac{g(P_{k-1})-g(P_k)}{P_{k-1}-P_k}$ if $P_{k-1}>P_k$, and
$N_k=N_{k-1}$ if $P_{k-1}=P_k$. Let $N_d=g(P_{d-1})/P_{d-1}$ and $\nu_k=N_{k+1}-N_k$. Write
$\operatorname{sl}(x,y)=\frac{g(y)-g(x)}{y-x}$ for $0\le x<y$; concavity gives
$\operatorname{sl}(y,z)\le\operatorname{sl}(x,y)$ for $0\le x<y<z$.

\emph{Claim 1: for $k\in[d]$ and $0\le x<P_k$, $N_k\le\operatorname{sl}(x,P_k)$.} For $k=0$ this holds
because $g$ is nondecreasing. Assume the claim for $k$, and let $k+1\le d-1$ and $0\le x<P_{k+1}$. If
$P_{k+1}=P_k$ then
$N_{k+1}=N_k\le\operatorname{sl}(x,P_k)=\operatorname{sl}(x,P_{k+1})$. Otherwise
$N_{k+1}=\operatorname{sl}(P_{k+1},P_k)\le\operatorname{sl}(x,P_{k+1})$ by concavity.

\emph{Claim 2: $\nu_k\ge0$.} For $k\le d-2$: if $P_{k+1}=P_k$ then $\nu_k=0$; otherwise Claim~1 with
$x=P_{k+1}$ gives $N_k\le\operatorname{sl}(P_{k+1},P_k)=N_{k+1}$. For $k=d-1$, Claim~1 with $x=0$ gives
$N_{d-1}\le(g(P_{d-1})-g(0))/P_{d-1}\le N_d$, since $g(0)\ge0$.

\emph{Claim 3: the identity.} In all cases $N_{k+1}(P_k-P_{k+1})=g(P_k)-g(P_{k+1})$ for $k\le d-2$, and
$N_dP_{d-1}=g(P_{d-1})$. Since $\min(P_i,P_k)=P_i$ for $k<i$ and $=P_k$ for $k\ge i$, and
$\sum_{k<i}\nu_k=N_i$,
\begin{align*}
  \sum_k\nu_k\min(P_i,P_k)&=N_iP_i+\sum_{k=i}^{d-1}(N_{k+1}-N_k)P_k\\
  &=\sum_{k=i}^{d-2}N_{k+1}(P_k-P_{k+1})+N_dP_{d-1}=g(P_i) .\qedhere
\end{align*}
\end{proof}

\begin{corollary}[thermal targets]\label{cor:thermal}
Let $E_0\le\dots\le E_{d-1}$, $\beta>0$, $0\le\beta'\le\beta$, $p=\gamma_\beta(E)$ and $q=\gamma_{\beta'}(E)$.
For $k\in[d]$ put
\[
  Z_k=\sum_{i}\min(p_i,p_k)>0,\qquad w_k=\min(p,p_k)/Z_k ,
\]
the \emph{clipped points} of $p$. Then $q=\sum_k\mu_kw_k$ with $\mu_k\ge0$ and $\sum_k\mu_k=1$.
\end{corollary}

\begin{proof}
$p$ is positive and nonincreasing. With $t=\beta'/\beta\in[0,1]$ we have
$p_i^t=e^{-\beta'E_i}/(\sum_je^{-\beta E_j})^t$, hence $q_i=p_i^t/S$ with $S=\sum_jp_j^t$. The function
$g(x)=x^t$ (with $g\equiv1$ if $t=0$) is concave and nondecreasing on $[0,\infty)$ with $g(0)\ge0$.
Lemma~\ref{lem:T} gives $\nu_k\ge0$ with $p_i^t=\sum_k\nu_k\min(p_i,p_k)$. Put $\mu_k=\nu_kZ_k/S\ge0$.
Then $\sum_k\mu_kw_k=q$, and $\sum_k\mu_k=\sum_k\nu_k\sum_i\min(p_i,p_k)/S=\sum_ip_i^t/S=1$.
\end{proof}

In Lean these are \lean{lemmaT} and \lean{thermal\_clip\_weights} (\lean{KirwanFreeCore.lean}).
By Lemma~\ref{lem:convex} and Proposition~\ref{prop:realisation}, Theorem~\ref{thm:diag} follows once
every clipped point $w_k$ lies in $\Reach(p)$. This is the content of Proposition~\ref{prop:wk} below.

\subsection{The decomposition of \texorpdfstring{$p\otimes p$}{p x p}}
Fix $k\in[d]$. Put $p_d=0$ and
\[
  c_m=p_m-p_{m+1}\ \ (m\in[d]),\qquad B_i=\min(p_i,p_k),\qquad Z=\sum_iB_i=Z_k .
\]
Then $c_m\ge0$ and $c_{d-1}=p_{d-1}>0$.

\begin{lemma}\label{lem:telescope}
For all $i\in[d]$:
\begin{align*}
&\text{(I1)}\quad p_i=\sum_{M=0}^{d-1}\ind{i\le M}c_M, &&\text{(I2)}\quad B_i=\sum_{M=k}^{d-1}\ind{i\le M}c_M,\\
&\text{(I3)}\quad p_i-B_i=\sum_{m=0}^{k-1}\ind{i\le m}c_m, &&\text{(I4)}\quad Z=\sum_{M=k}^{d-1}(M+1)c_M,\\
&\text{(I5)}\quad 1-Z=\sum_{m=0}^{k-1}(m+1)c_m .
\end{align*}
\end{lemma}

\begin{proof}
(I1) telescopes. For (I2), the right-hand side is $p_{\max(i,k)}=\min(p_i,p_k)$ because $p$ is
nonincreasing. (I3) is (I1)$-$(I2). Summing (I2) and (I3) over $i$ gives (I4) and (I5), using
$\sum_ip_i=1$.
\end{proof}

For $u,v\in\N$ let $C(u,v)=([0,u]\times[0,v])\cup([0,v]\times[0,u])\subseteq\N^2$ be the \emph{cross}
with arms $u,v$, and $\chi_{u,v}$ its indicator function. It is symmetric in $(u,v)$ and in its two
arguments, and $Q_u=\chi_{u,u}$ is the indicator of the square $[0,u]^2$. Crosses and squares are
symmetric Young diagrams (down-closed subsets of $\N^2$). (In the working notes accompanying the Lean
code the squares are called $Q$ and the crosses $U$; in Lean, $\chi_{u,v}$ is \lean{crossInd u v}.)

\begin{lemma}[decomposition]\label{lem:decomp}
Let $\omega_{m,M}=\tfrac12c_mc_M(1+\ind{M\ge k})$ for $m<k$ and $M\in[d]$. For all $i,j\in[d]$,
\[
  p_ip_j=B_iB_j+\sum_{m=0}^{k-1}\sum_{M=0}^{d-1}\omega_{m,M}\bigl(Q_{\min(m,M)}(i,j)+\chi_{m,M}(i,j)\bigr) .
\]
\end{lemma}

\begin{proof}
By inclusion--exclusion,
$Q_{\min(m,M)}(i,j)+\chi_{m,M}(i,j)=\ind{i\le m}\ind{j\le M}+\ind{i\le M}\ind{j\le m}$. By (I1) and (I2),
$\sum_M\tfrac12c_M(1+\ind{M\ge k})\ind{j\le M}=\tfrac12(p_j+B_j)$. With (I3) the double sum becomes
$\tfrac12(p_i-B_i)(p_j+B_j)+\tfrac12(p_j-B_j)(p_i+B_i)=p_ip_j-B_iB_j$.
\end{proof}

\begin{lemma}[runs of a cross]\label{lem:crossrun}
Let $u\le v$ and $s\ge0$. For $i\in\N$, $\chi_{u,v}(i,i+s)=1$ if and only if $s\le v$ and
$i\le\min(u,v-s)$. Consequently, for $v<d$ the run $(\chi_{u,v}(i,i+s))_{i<d-s}$ equals
$\mathbf 1_{N_s}$ with $N_s=\#\{a\in[v-u,v]:a\ge s\}$, and $C(u,v)\cap[d]^2$ has
$\ms(v-u,v)$ cells, where
\[
  \ms(\ell,h)=\sum_{a=\ell}^{h}(2a+1)=(h+1)^2-\ell^2 .
\]
\end{lemma}

\begin{proof}
The condition $(i\le v\wedge i+s\le u)$ implies $(i\le u\wedge i+s\le v)$ because $u\le v$, so
$\chi_{u,v}(i,i+s)=1$ iff $i\le u$ and $i+s\le v$. For $s\le v$ this gives
$\min(u,v-s)+1=\#\{a\in[v-u,v]:a\ge s\}$ ones at the start of the run; for $s>v$ both sides are $0$.
The cell count is $(u+1)(v+1)+(v+1)(u+1)-(u+1)^2=(v+1)^2-(v-u)^2$.
\end{proof}

Thus $C(u,v)$ is the symmetric Young diagram whose Frobenius hooks have the arm lengths
$v-u,\dots,v$: its runs only depend on the \emph{arm interval} $[v-u,v]$. The square $Q_u$ has arm
interval $[0,u]$.

\subsection{The Hall identity}
For $k\le M\le d-1$ define the \emph{capacity} of the box $[0,M]$, the mass of the \emph{big crosses}
assigned to it, and the \emph{leftover}:
\[
  A_M=\frac{(1-Z^2)(M+1)c_M}{Z},\qquad \Lambda_M=\sum_{m=0}^{k-1}\omega_{m,M}\,\ms(M-m,M),\qquad
  L_M=A_M-\Lambda_M,
\]
and $L=\sum_{M=k}^{d-1}L_M$.

\begin{lemma}[Hall identity]\label{lem:hall}
For $k\le M\le d-1$,
\[
  L_M=c_M\Bigl(\frac{(M+1)(1-Z)^2}{Z}+S_k\Bigr),\qquad S_k=\sum_{i=0}^{k-1}(2i+1)(p_i-p_k)\ge0 .
\]
In particular $L_M\ge0$. Moreover $\sum_{M=k}^{d-1}A_M=1-Z^2$, and $L>0$ unless $p_0=p_1=\dots=p_k$.
\end{lemma}

\begin{proof}
For $M\ge k$ we have $\omega_{m,M}=c_mc_M$ and $\ms(M-m,M)=(m+1)(2M+1-m)=2(m+1)(M+1)-(m+1)^2$. By (I5)
and the summation by parts $\sum_{m<k}(m+1)^2c_m=\sum_{m<k}\sum_{i\le m}(2i+1)c_m=S_k$, we get
$\Lambda_M=c_M(2(M+1)(1-Z)-S_k)$. Hence
\[
  L_M=c_M\Bigl((M+1)(1-Z)\bigl(\tfrac{1+Z}{Z}-2\bigr)+S_k\Bigr)=c_M\Bigl(\tfrac{(M+1)(1-Z)^2}{Z}+S_k\Bigr).
\]
By (I4), $\sum_{M\ge k}A_M=(1-Z^2)Z/Z$. If $p_m>p_{m+1}$ for some $m<k$, then $1-Z\ge(m+1)c_m>0$ by
(I5), so $L\ge L_{d-1}\ge d\,c_{d-1}(1-Z)^2/Z>0$.
\end{proof}

The Lean development proves the inequality $L_M\ge(M+1)c_M(1-Z)^2/Z$, which suffices
(\lean{leftM\_nonneg}, \lean{Ltot\_pos}, \lean{sum\_aM} in \lean{KirwanFree.lean}, using the bound
$\ms(M-m,M)\le2(m+1)(M+1)$); Lemma~\ref{lem:decomp} is \lean{wk\_cell\_identity} there.

%==================================================================================================
\section{Interval uniformisations}\label{sec:unif}
%==================================================================================================

\subsection{Definition}
The following definition is, clause by clause, the Lean structure \lean{IsIntervalUnif n lo hi u}
(\lean{IntervalUnif.lean}), with $\ell=$ \lean{lo} and $h=$ \lean{hi}.

\begin{definition}[interval uniformisation]\label{def:unif}
Let $n\ge1$ and $0\le\ell\le h\le n-1$. An \emph{interval uniformisation of $[\ell,h]$ in the $n$-box} is
a function $Y:\N\times\N\to\R$ such that
\begin{enumerate}
\item[(U1)] $Y(i,j)=Y(j,i)$ for all $i,j$;
\item[(U2)] $0\le Y(i,j)\le1$ for all $i,j$;
\item[(U3)] $Y(i,j)=0$ whenever $i\ge n$;
\item[(U4)] $\sum_{i=0}^{n-1}Y(i,i+s)=r_s:=\#\{a\in[\ell,h]:a\ge s\}$ for every $s\ge0$;
\item[(U5)] $\sum_{j=0}^{n-1}Y(r,j)=\frac1n\sum_{a=\ell}^{h}(2a+1)=\frac{\ms(\ell,h)}{n}$ for every $r\in[n]$.
\end{enumerate}
We say that $\mathrm{UNIF}(n)$ holds (Lean: \lean{IntervalUnifAll n}) if such a $Y$ exists for every
$0\le\ell\le h\le n-1$.
\end{definition}

By (U1) and (U3), $Y$ vanishes outside $[n]^2$. The runs of $Y$ have the totals of the runs of a
symmetric Young diagram with arm interval $[\ell,h]$, but all row sums are equal.

\begin{lemma}\label{lem:unifdom}
Let $Y$ be an interval uniformisation of $[v-u,v]$ in the $n$-box, where $u\le v<n\le d$. Then for every
$s<d$ the run $(Y(i,i+s))_{i<d-s}$ satisfies $(Y(i,i+s))_{i<d-s}\dom(\chi_{u,v}(i,i+s))_{i<d-s}$, and
for $r\in[d]$, $\sum_{j<d}Y(r,j)=\ind{r<n}\,\ms(v-u,v)/n$.
\end{lemma}

\begin{proof}
The entries lie in $[0,1]$ by (U2). Since $Y(i,i+s)=0$ when $i+s\ge n$ (by (U1), (U3)) and $n\le d$,
both $\sum_{i<d-s}Y(i,i+s)$ and the sum in (U4) reduce to $\sum_{i<n-s}Y(i,i+s)$, so the run sums to
$r_s$. By Lemma~\ref{lem:crossrun} the cross run is $\mathbf 1_{r_s}$, and Lemma~\ref{lem:dom}(d)
applies. The row sums follow from (U3) and (U5).
\end{proof}

In Lean this is \lean{IsIntervalUnif.smaj} and \lean{IsIntervalUnif.row'}.

\subsection{Clipped points are reachable}

\begin{proposition}\label{prop:wk}
Assume that for every $n\le d$ and every $0\le\ell\le h\le n-1$ an interval uniformisation
$Y_n[\ell,h]$ of $[\ell,h]$ in the $n$-box is given. Then $w_k=B/Z\in\Reach(p)$ for every $k\in[d]$.
\end{proposition}

\begin{proof}
Fix $k$ and use the notation of Section~\ref{sec:reduction}. All matrices below are restricted to
$[d]\times[d]$.

\emph{Trivial case.} If $p_0=\dots=p_k$, then $B=p$, $Z=1$ and $w_k=p$; the configuration
$F(i,j)=p_ip_j$ is admissible by Lemma~\ref{lem:dom}(c) and has row sums $p$.

\emph{General case.} Otherwise $L>0$ by Lemma~\ref{lem:hall}. Put $\xi_M=L_M/L$ for $k\le M\le d-1$,
so $\xi_M\ge0$ and $\sum_M\xi_M=1$. For $0\le\ell\le h\le k-1$ define the \emph{flexible piece}
\[
  Y^{\mathrm{fl}}[\ell,h]=\sum_{M=k}^{d-1}\xi_M\,Y_{M+1}[\ell,h] ,
\]
a convex combination of uniformisations in the boxes $[0,M]$, $M\ge k>h$. For $m<k$ and $M\in[d]$ let
\[
  G_{m,M}=
  \begin{cases}
    Y_{M+1}[M-m,M]+Y^{\mathrm{fl}}[0,m], & M\ge k,\\[2pt]
    Y^{\mathrm{fl}}[\,|m-M|,\max(m,M)]+Y^{\mathrm{fl}}[0,\min(m,M)], & M<k,
  \end{cases}
\]
and
\[
  F=B\otimes B+\sum_{m=0}^{k-1}\sum_{M=0}^{d-1}\omega_{m,M}\,G_{m,M} .
\]
Thus the square $Q_{\min(m,M)}$ is always replaced by a flexible piece; a cross $\chi_{m,M}$ with
$M\ge k$ (a \emph{big cross}) is replaced by a uniformisation of its arm interval $[M-m,M]$ in its own
box $[0,M]$ (then $\min(m,M)=m$); and a cross with $M<k$ is replaced by a flexible piece.

\emph{Symmetry.} Every term is symmetric.

\emph{Runs.} Fix $s<d$. The run of $B\otimes B$ is nonincreasing, so it is $\dom$ itself
(Lemma~\ref{lem:dom}(c)). By Lemma~\ref{lem:unifdom} (with $n=M+1\le d$), the run of
$Y_{M+1}[M-m,M]$ is $\dom$ the run of $\chi_{m,M}$, and for $v<k\le M$ the run of
$Y_{M+1}[v-u,v]$ is $\dom$ the run of $\chi_{u,v}$; averaging with the weights $\xi_M$
(Lemma~\ref{lem:dom}(b)), the run of $Y^{\mathrm{fl}}[v-u,v]$ is $\dom$ the run of $\chi_{u,v}$. Hence
the run of $G_{m,M}$ is $\dom$ the run of $Q_{\min(m,M)}+\chi_{m,M}$. Adding with the weights
$\omega_{m,M}\ge0$ and using Lemma~\ref{lem:decomp}, $F^{(s)}\dom\lambda^{(s)}$.

\emph{Rows.} Fix $r\in[d]$ and let $\theta(r)=\sum_{M=k}^{d-1}\xi_M\ind{r\le M}/(M+1)$. By
Lemma~\ref{lem:unifdom}, row $r$ of $Y_{M+1}[\ell,h]$ sums to $\ms(\ell,h)\ind{r\le M}/(M+1)$ and row
$r$ of $Y^{\mathrm{fl}}[\ell,h]$ to $\ms(\ell,h)\theta(r)$. Row $r$ of $B\otimes B$ sums to $B_rZ$.
Collecting terms,
\[
  \sum_jF(r,j)=B_rZ+\sum_{M=k}^{d-1}\Lambda_M\frac{\ind{r\le M}}{M+1}+\Lambda^{\mathrm{fl}}\,\theta(r),
\]
where $\Lambda^{\mathrm{fl}}$ is the total $\omega$-weighted mass of the flexible pieces. The total mass
of all pieces is
$\sum_{m,M}\omega_{m,M}\bigl(\ms(0,\min(m,M))+\ms(|m-M|,\max(m,M))\bigr)=\sum_{i,j<d}(p_ip_j-B_iB_j)=1-Z^2$ by
Lemmas~\ref{lem:decomp} and~\ref{lem:crossrun}, so
$\Lambda^{\mathrm{fl}}=1-Z^2-\sum_M\Lambda_M=\sum_M(A_M-\Lambda_M)=L$ by Lemma~\ref{lem:hall}. Since
$\xi_ML=L_M$,
\begin{align*}
  \sum_jF(r,j)&=B_rZ+\sum_{M=k}^{d-1}(\Lambda_M+L_M)\frac{\ind{r\le M}}{M+1}
  =B_rZ+\frac{1-Z^2}{Z}\sum_{M=k}^{d-1}\ind{r\le M}c_M\\
  &=B_rZ+\frac{(1-Z^2)B_r}{Z}=\frac{B_r}{Z},
\end{align*}
using the definition of $A_M$ and (I2). Hence $F$ is admissible with row sums $w_k$.
\end{proof}

In Lean this is \lean{wkU\_reach} (with \lean{wkU\_smaj} and \lean{wkU\_rows}) in
\lean{KirwanFreeU.lean}. The boxes used have size $n=M+1\le d$; the big crosses need the intervals
$[M-m,M]$ with $h=n-1$, and the flexible pieces need intervals with $h\le n-2$.

\subsection{The reduction and the main theorem}

\begin{proposition}[reduction]\label{prop:reduction}
If $\mathrm{UNIF}(n)$ holds for every $n\le d$, then Theorem~\ref{thm:diag} holds in dimension $d$.
\end{proposition}

\begin{proof}
For $d=0$ there is nothing to prove. For $d\ge1$ let $p=\gamma_\beta(E)$ and $q=\gamma_{\beta'}(E)$. By
Corollary~\ref{cor:thermal}, $q$ is a convex combination of the $w_k$, which lie in $\Reach(p)$ by
Proposition~\ref{prop:wk}. By Lemma~\ref{lem:convex}, $q\in\Reach(p)$, and
Proposition~\ref{prop:realisation} gives a real orthogonal $O$ with both marginals of
$O\diag(p\otimes p)O^{\mathsf T}$ equal to $\diag(q)$. Since $D_\beta\otimes D_\beta=\diag(p\otimes p)$
and $D_{\beta'}=\diag(q)$, the matrix $U=O$ satisfies Theorem~\ref{thm:diag}.
\end{proof}

This is \lean{stu\_exists\_of\_intervalUnif} (\lean{KirwanFreeU.lean}). The remaining ingredient is:

\begin{mainthm}[interval uniformisations exist]\label{thm:unif}
$\mathrm{UNIF}(n)$ holds for every $n\ge1$.
\end{mainthm}

\begin{proof}[Proof of Theorem~\ref{thm:diag}]
Combine Proposition~\ref{prop:reduction} and Theorem~\ref{thm:unif}. (In Lean:
\lean{stu\_exists\_unconditional}, in \lean{KirwanFreeFinal.lean}.)
\end{proof}

Theorem~\ref{thm:unif} is proved in Section~\ref{sec:proofunif}. The constructions are summarised in
Table~\ref{tab:cases}, where $K=n-1-h$ and $M=n-\ell$.

\begin{table}[ht]
\centering
\begin{tabular}{llll}
\toprule
interval & case & construction & reference\\
\midrule
$[0,n-1]$ & (A) $h=n-1$ & $J_n$ (all ones) & Lemma~\ref{lem:compl}\\
$[\ell,n-1]$, $\ell\ge1$ & (A) $h=n-1$ & $J_n-\cR^{\ell-1}_n$ & Lemma~\ref{lem:compl}\\
$[0,h]$, $h\le n-2$ & (B) $\ell=0$ & renewal square $\cR^h_n$ & Proposition~\ref{prop:renewal}\\
$[\ell,h]$, $1\le\ell\le h\le n-2$ & (C) $n\le2K+1$ & $Y^{0}$ (no dipole) & Proposition~\ref{prop:caseC}\\
same & (D) $K=1$, $n\ge4$ & $Y^{2/7}$ & Proposition~\ref{prop:caseD}\\
same & (E) $K\ge2$, $n\ge2K+2$ & $Y^{\alpha^*}$ & Proposition~\ref{prop:caseE}\\
\bottomrule
\end{tabular}
\caption{The interval uniformisations. $Y^\alpha$ is the T/H/G matrix with one dipole of amplitude
$\alpha$ (Definition~\ref{def:dipole}); $\alpha^*$ is given in \eqref{eq:alphastar}.}
\label{tab:cases}
\end{table}

\subsection{Renewal squares and complements}

\begin{proposition}[renewal squares]\label{prop:renewal}
Let $n\ge2$, $0\le m\le n-2$ and $K=n-m-1\ge1$. Define $\eta_\ell=1-\frac{K^2}{(n-\ell)(n-\ell+1)}$ for
$1\le\ell\le m+1$ and $\eta_\ell=0$ for $\ell=0$ and for $\ell\ge m+2$; put
$A_r=\sum_{t=0}^{m+1}(-1)^t\eta_{r+t}$ for $r\ge0$, and
\[
  \cR^m_n(i,j)=A_{|i-j|+1}+A_{\min(i+j+2,\,2n-i-j)}\quad(i,j<n),\qquad \cR^m_n(i,j)=0\ \text{otherwise}.
\]
Then $\cR^m_n$ is an interval uniformisation of $[0,m]$ in the $n$-box.
\end{proposition}

\begin{proof}
\emph{(i) Tails.} In the sum $A_r+A_{r+1}$ all terms cancel in pairs except the first term of $A_r$
and the last term of $A_{r+1}$, so $A_r+A_{r+1}=\eta_r+(-1)^{m+1}\eta_{r+m+2}=\eta_r$, because
$r+m+2\ge m+2$.

\emph{(ii) Weights.} For $1\le\ell\le m$,
$\eta_\ell-\eta_{\ell+1}=\frac{2K^2}{(n-\ell-1)(n-\ell)(n-\ell+1)}\ge0$, and $\eta_{m+1}=\frac1{K+1}$. Since
$(n-\ell)(n-\ell+1)\ge K(K+1)>K^2$ for $\ell\le m+1$, we get
$1\ge\eta_1\ge\eta_2\ge\dots\ge\eta_{m+1}>0=\eta_{m+2}=\eta_{m+3}=\cdots$.

\emph{(iii) Alternating sums.} For $r\ge1$ and $\kappa\ge0$, (i) gives
$A_r+A_{r+2\kappa+1}=\sum_{t=0}^{2\kappa}(-1)^t\eta_{r+t}$, an alternating sum of a nonincreasing
nonnegative sequence; hence $0\le A_r+A_{r+2\kappa+1}\le\eta_r\le1$ (by induction on $\kappa$:
$A_r+A_{r+2\kappa+3}=(\eta_r-\eta_{r+1})+(A_{r+2}+A_{r+2\kappa+3})$).

\emph{(iv) Entries.} For $i,j<n$ one checks $\min(i+j+2,2n-i-j)=|i-j|+2+2\kappa$ with
$\kappa=\min(i,j,n-1-i,n-1-j)$. So $\cR^m_n(i,j)=A_{s+1}+A_{s+1+2\kappa+1}\in[0,\eta_{s+1}]\subseteq[0,1]$
with $s=|i-j|$, and $\cR^m_n(i,j)=0$ if $s\ge m+1$, since then both indices are $\ge m+2$ and
$A_r=0$ for $r\ge m+2$. (U1)--(U3) hold.

\emph{(v) Rows.} Row $0$: $\cR^m_n(0,j)=A_{j+1}+A_{j+2}=\eta_{j+1}$, so it sums to
$\sum_{\ell=1}^{m+1}\eta_\ell=(m+1)-K^2\bigl(\frac1{n-m-1}-\frac1n\bigr)=\frac{(m+1)^2}{n}=\frac{\ms(0,m)}n$.
For $0\le r\le n-2$, substituting $j\mapsto j+1$ shows that the Toeplitz part changes by
$\sum_jA_{|r+1-j|+1}-\sum_jA_{|r-j|+1}=A_{r+2}-A_{n-r}$, and the Hankel part by
$\sum_jA_{\varphi(r+1+j)}-\sum_jA_{\varphi(r+j)}=A_{\varphi(r+n)}-A_{\varphi(r)}=A_{n-r}-A_{r+2}$,
where $\varphi(S)=\min(S+2,2n-S)$ (all sums over $0\le j\le n-1$). So all rows have the same sum. This is (U5).

\emph{(vi) Runs.} Let $W_\ell=2\eta_\ell+(n-\ell-1)(\eta_\ell-\eta_{\ell+1})$. By (ii), $W_\ell=2$ for
$1\le\ell\le m$ and $W_{m+1}=(n-m)\eta_{m+1}=1$. For $0\le s\le n-2$ we claim
$\run_s=W_{s+1}+\run_{s+2}$, where $\run_s=\sum_i\cR^m_n(i,i+s)$. The first and last cells of run $s$,
$(0,s)$ and $(n-1-s,n-1)$, have $\kappa=0$ and value $A_{s+1}+A_{s+2}=\eta_{s+1}$. For
$0\le i\le n-s-3$, the cells $(i+1,i+1+s)$ and $(i,i+s+2)$ have the same $i+j$, hence the same Hankel
term, and Toeplitz terms $A_{s+1}$ and $A_{s+3}$, which differ by $\eta_{s+1}-\eta_{s+2}$ by (i). The
cells $(i,i+s+2)$, $0\le i\le n-s-3$, form run $s+2$. Summing,
$\run_s=2\eta_{s+1}+\run_{s+2}+(n-s-2)(\eta_{s+1}-\eta_{s+2})=W_{s+1}+\run_{s+2}$. Now $\run_s=0$ for
$s\ge m+1$ by (iv), $\run_m=W_{m+1}=1$, $\run_{m-1}=W_m=2$ if $m\ge1$, and downward induction gives
$\run_s=m+1-s=\#\{a\in[0,m]:a\ge s\}$ for $s\le m$. This is (U4).
\end{proof}

The name comes from an equivalent description: $\cR^m_n$ is the superposition of the reversal
permutation matrices of the blocks $[0,\ell-1]$ and $[n-\ell,n-1]$ with weight $\eta_\ell$ and of the
interior blocks $[a,a+\ell-1]$ ($1\le a\le n-1-\ell$) with weight $\eta_\ell-\eta_{\ell+1}$, a
stationary renewal process of block boundaries. We do not need this description (it was checked for
$n\le24$, Appendix~\ref{app:checks}). In Lean, $\cR^m_n$ is \lean{renU} and the proposition is
\lean{isIntervalUnif\_renU} (\lean{Renewal.lean}).

\begin{lemma}[$J$ and complements]\label{lem:compl}
The all-ones matrix $J_n$ of the $n$-box is an interval uniformisation of $[0,n-1]$. If $Y$ is an
interval uniformisation of $[0,m]$ in the $n$-box, $m\le n-2$, then $J_n-Y$ (restricted to the box) is
one of $[m+1,n-1]$.
\end{lemma}

\begin{proof}
For $J_n$: $\run_s=n-s=\#\{a\in[0,n-1]:a\ge s\}$ and all rows equal $n=n^2/n$. For $J_n-Y$: entries lie in
$[0,1]$; $\run_s=(n-s)-\#\{a\in[0,m]:a\ge s\}=\#\{a\in[m+1,n-1]:a\ge s\}$; rows equal
$n-(m+1)^2/n=\ms(m+1,n-1)/n$.
\end{proof}

In Lean: \lean{isIntervalUnif\_boxJ} and \lean{IsIntervalUnif.compl} (\lean{Renewal.lean}).

%==================================================================================================
\section{T/H/G matrices}\label{sec:thg}
%==================================================================================================

To avoid confusion with the Hamiltonian, the three sequences below are written in sans-serif.

\begin{definition}\label{def:thg}
Let $n\ge1$ and $\sT,\sH,\sG:\N\to\R$. The \emph{T/H/G matrix} $Y=Y[\sT,\sH,\sG]$ of the $n$-box is
defined for $i,j<n$ by
\[
  Y(i,j)=
  \begin{cases}
    \tfrac12\bigl(\sT(|i-j|)+\sH(m'(i,j))\bigr), & i+j\ne n-1,\\[2pt]
    \tfrac12\bigl(\sT(|i-j|)+\sG(|i-j|)\bigr), & i+j=n-1,
  \end{cases}
  \qquad m'(i,j)=\min(i+j+1,\,2n-1-i-j),
\]
and $Y(i,j)=0$ if $i\ge n$ or $j\ge n$. We write $\run_s=\sum_{i<n}Y(i,i+s)$ and
$\row_r=\sum_{j<n}Y(r,j)$.
\end{definition}

The Toeplitz part depends on $|i-j|$, the Hankel part on $i+j$ (through the ``level'' $m'$), and the
antidiagonal $i+j=n-1$ is treated separately.

\begin{lemma}[levels and cells]\label{lem:levels}
Let $i,j<n$ and $s=|i-j|$. Then $m'(i,j)-s$ is odd and positive, $m'(i,j)\le n$, and $m'(i,j)=n$ if
and only if $i+j=n-1$; if $i+j=n-1$ then $n-s$ is odd. Consequently $Y$ is symmetric, and $0\le Y\le1$
holds provided
\begin{align}
  0\le\sT(s)+\sH(u)\le2&\quad\text{for all }0\le s<u\le n-1\text{ with }u-s\text{ odd},\label{eq:cellTH}\\
  0\le\sT(s)+\sG(s)\le2&\quad\text{for all }0\le s\le n-1\text{ with }n-s\text{ odd}.\label{eq:cellTG}
\end{align}
\end{lemma}

\begin{proof}
$i+j+1-s=2\min(i,j)+1$ and $2n-1-i-j-s=2n-1-2\max(i,j)\ge1$ are odd and positive. If $i+j\le n-2$ then
$m'\le i+j+1\le n-1$; if $i+j\ge n$ then $m'\le2n-1-i-j\le n-1$; if $i+j=n-1$ then $m'=n$ and
$n-s=2\min(i,j)+1$.
\end{proof}

\begin{lemma}[runs]\label{lem:thgruns}
\textup{(a)} $\run_s=0$ for $s\ge n$. \textup{(b)} $\run_{n-1}=\frac12(\sT(n-1)+\sG(n-1))$.
\textup{(c)} If $n\ge2$, $\run_{n-2}=\sT(n-2)+\sH(n-1)$. \textup{(d)} For $0\le s\le n-3$,
\begin{multline*}
  2\run_s-2\run_{s+2}\\
  =(n-s)\sT(s)-(n-s-2)\sT(s+2)+2\sH(s+1)+\ind{n-s\text{ odd}}\bigl(\sG(s)-\sG(s+2)\bigr).
\end{multline*}
\end{lemma}

\begin{proof}
(a) is clear. (b) Run $n-1$ is the antidiagonal cell $(0,n-1)$. (c) Run $n-2$ consists of $(0,n-2)$
and $(1,n-1)$, with $i+j\in\{n-2,n\}$ and $m'=n-1$ in both cases. (d) Let $0\le i\le n-s-3$. The cells
$(i+1,i+1+s)$ and $(i,i+s+2)$ have the same sum $i+j=2i+s+2$, so both are antidiagonal or neither is, and
they have the same level $m'$. Their difference is $\frac12(\sT(s)-\sT(s+2))$, plus
$\frac12(\sG(s)-\sG(s+2))$ when $2i+s+3=n$; such an $i$ exists (and is unique) iff $n-s$ is odd. The
cells $(i,i+s+2)$, $0\le i\le n-s-3$, form run $s+2$. The two remaining cells of run $s$, $(0,s)$ and
$(n-1-s,n-1)$, are not antidiagonal and have level $s+1$, so each equals $\frac12(\sT(s)+\sH(s+1))$.
Summing and multiplying by $2$ gives the claim.
\end{proof}

\begin{lemma}[row shift]\label{lem:thgrows}
For $0\le r\le n-2$,
\begin{multline*}
  2\row_{r+1}-2\row_r=\bigl(\sT(r+1)-\sH(r+1)\bigr)-\bigl(\sT(n-1-r)-\sH(n-1-r)\bigr)\\
  +\sG(|2r+3-n|)-\sG(|2r+1-n|) .
\end{multline*}
\end{lemma}

\begin{proof}
Let $\sH^*(S)=\sH(\min(S+1,2n-1-S))$ for $S\ne n-1$ and $\sH^*(n-1)=0$. The antidiagonal cell of row $r$
is $j=n-1-r$, at distance $|2r+1-n|$, so
$2\row_r=\sum_{j<n}\sT(|r-j|)+\sum_{j<n}\sH^*(r+j)+\sG(|2r+1-n|)$. Shifting $r$ to $r+1$, the Toeplitz sum
gains $\sT(r+1)$ (the term $j=0$) and loses $\sT(n-1-r)$ (the old term $j=n-1$); the Hankel sum
$\sum_{S=r+1}^{r+n}\sH^*(S)$ loses $\sH^*(r)=\sH(r+1)$ and gains $\sH^*(r+n)=\sH(n-1-r)$; the antidiagonal
term changes from $\sG(|2r+1-n|)$ to $\sG(|2r+3-n|)$.
\end{proof}

\begin{lemma}[the row value is forced]\label{lem:rowvalue}
Let $Y:\N\times\N\to\R$ satisfy (U1), (U3) and (U4) for $[\ell,h]$, $h<n$, and suppose all rows
$r\in[n]$ have the same sum $c$. Then $c=\ms(\ell,h)/n$, i.e.\ (U5) holds.
\end{lemma}

\begin{proof}
Splitting $[n]^2$ into the upper runs and the strictly lower runs and using symmetry,
$nc=\sum_{s\ge0}\run_s+\sum_{s\ge1}\run_s=\sum_{s\ge0}(r_s+r_{s+1})=\sum_{a=\ell}^{h}\bigl((a+1)+a\bigr)$.
\end{proof}

In Lean these are \lean{thgU}, \lean{thg\_cells}, \lean{thg\_run\_rec}, \lean{thg\_run\_top},
\lean{thg\_run\_top2}, \lean{thg\_row\_succ} and \lean{row\_const\_value} (\lean{DipoleUnif.lean}).

%==================================================================================================
\section{The one-dipole construction}\label{sec:dipole}
%==================================================================================================

\begin{definition}\label{def:dipole}
Let $1\le\ell\le h\le n-2$ and $\alpha\in\R$, and put $K=n-1-h$ and $M=n-\ell$, so that $K\ge1$ and
$K+1\le M\le n-1$. For integers $x\ge0$ and $L\ge0$ let
\[
  q_L(x)=\ind{x\ge L+1}\bigl(2x(x-1)-2L^2\bigr),\qquad
  \cQ(x)=q_K(x)-q_M(x)+2(K+1)\alpha\,\ind{K+2\le x\le2K+2},
\]
and define $\vartheta:\N\to\R$ by $\vartheta(0)=\vartheta(1)=0$ and
\[
  \vartheta(x)=\frac{\cQ(x)}{x(x-1)}-\vartheta(x-1)\qquad(x\ge2).
\]
Set $\sT(s)=\vartheta(n-s)$ for $0\le s\le n$ and $\sT(s)=0$ for $s>n$,
\[
  \sH(u)=\sT(u)-\alpha\ind{u=h},\qquad \sG(u)=\alpha\ind{u\le2h-n-1},
\]
and let $Y^\alpha$ be the T/H/G matrix $Y[\sT,\sH,\sG]$ of the $n$-box.
\end{definition}

Here $x=n-s$ is the natural coordinate: $\vartheta$ depends only on $K$, $M$ and $\alpha$, not on $n$,
which only truncates the range to $x\le n$. The pair ($-\alpha$ at the Hankel level $h$, $+\alpha$ on the
antidiagonal cells with $|i-j|\le2h-n-1$) is the \emph{dipole}; for $\alpha=0$ the construction is the
difference $\cR^h_n-\cR^{\ell-1}_n$ of two renewal squares (Remark~\ref{rem:parity}). In Lean,
$\cQ$, $\vartheta$ and $Y^\alpha$ are \lean{dipQ}, \lean{dipTau} and \lean{dipU}; Lean's definition
$\vartheta(x+1)=\cQ(x+1)/((x+1)x)-\vartheta(x)$ for all $x\ge0$ gives the same values, because Lean
divides by zero to zero at $x=0$.

\begin{lemma}\label{lem:Qdiff}
For every $x\ge0$,
\[
  \cQ(x+1)-\cQ(x)=x\Bigl(2\ind{x\ge K}+2\ind{x\ge K+1}-2\ind{x\ge M}-2\ind{x\ge M+1}
  +2\alpha\ind{x=K+1}-\alpha\ind{x=2K+2}\Bigr).
\]
\end{lemma}

\begin{proof}
$q_L(x+1)-q_L(x)$ equals $4x$ if $x\ge L+1$, $2L=2x$ if $x=L$, and $0$ if $x<L$; that is,
$x(2\ind{x\ge L}+2\ind{x\ge L+1})$. The window term changes by $+2(K+1)\alpha=2\alpha x$ at $x=K+1$ and
by $-2(K+1)\alpha=-\alpha x$ at $x=2K+2$.
\end{proof}

\begin{proposition}\label{prop:dipole}
In the setting of Definition~\ref{def:dipole}:
\begin{enumerate}
\item[(a)] for every $\alpha$, $\run_s(Y^\alpha)=r_s=\#\{a\in[\ell,h]:a\ge s\}$ for all $s\ge0$;
\item[(b)] if $\alpha=0$ or $2h\ge n$, all rows of $Y^\alpha$ have sum $\ms(\ell,h)/n$;
\item[(c)] if moreover
\begin{align}
  0\le\vartheta(x)+\vartheta(y)-\alpha\ind{y=K+1}\le2&\quad\text{for all }1\le y<x\le n\text{ with }x-y\text{ odd},
  \tag{CELL1}\label{eq:cell1}\\
  0\le\vartheta(x)+\alpha\ind{x\ge2K+3}\le2&\quad\text{for all odd }x\text{ with }1\le x\le n,
  \tag{CELL2}\label{eq:cell2}
\end{align}
then $Y^\alpha$ is an interval uniformisation of $[\ell,h]$ in the $n$-box.
\end{enumerate}
\end{proposition}

\begin{proof}
Note $n\ge3$. (a) By Lemma~\ref{lem:thgruns}(a), $\run_s=0=r_s$ for $s\ge n$. Next,
$\run_{n-1}=\frac12(\vartheta(1)+\alpha\ind{n-1\le2h-n-1})=0=r_{n-1}$, because $h\le n-2$. And
$\run_{n-2}=\vartheta(2)+\vartheta(1)-\alpha\ind{n-1=h}=\vartheta(2)=\cQ(2)/2=\ind{K=1}=\ind{h=n-2}=r_{n-2}$,
since $\cQ(2)=q_K(2)=2$ if $K=1$ and $\cQ(2)=0$ if $K\ge2$ (recall $M\ge2$). Since $r$ and $\run$ agree for
$s\in\{n-2,n-1\}$ and all $s\ge n$, it suffices by Lemma~\ref{lem:thgruns}(d) to show, for
$0\le s\le n-3$ and $y=n-s-2\ge1$,
\begin{equation}\label{eq:runstep}
  (n-s)\sT(s)-(n-s-2)\sT(s+2)+2\sH(s+1)+\ind{y\text{ odd}}\bigl(\sG(s)-\sG(s+2)\bigr)=2(r_s-r_{s+2}).
\end{equation}
Since $s+1=h$ iff $y=K$, the first three terms equal
\begin{align*}
  &(y+2)\bigl(\vartheta(y+2)+\vartheta(y+1)\bigr)-y\bigl(\vartheta(y+1)+\vartheta(y)\bigr)-2\alpha\ind{y=K}\\
  &\qquad=\frac{\cQ(y+2)-\cQ(y+1)}{y+1}-2\alpha\ind{y=K},
\end{align*}
by the definition of $\vartheta$ (at $y+2\ge3$ and $y+1\ge2$). By Lemma~\ref{lem:Qdiff} at $x=y+1$ this is
$2\ind{y\ge K-1}+2\ind{y\ge K}-2\ind{y\ge M-1}-2\ind{y\ge M}-\alpha\ind{y=2K+1}$. Next,
$\sG(s)-\sG(s+2)=\alpha\ind{s\in\{2h-n-2,2h-n-1\}}$, and $s=2h-n-1$ iff $y=2K+1$ (odd), while
$s=2h-n-2$ iff $y=2K+2$ (even); so the last term of \eqref{eq:runstep} is $\alpha\ind{y=2K+1}$. Hence the
left side of \eqref{eq:runstep} is $2\ind{K-1\le y\le M-2}+2\ind{K\le y\le M-1}$. On the right,
$r_s-r_{s+2}=\ind{\ell\le s\le h}+\ind{\ell\le s+1\le h}$, and $\ell\le s\le h$ iff $K-1\le y\le M-2$,
$\ell\le s+1\le h$ iff $K\le y\le M-1$.

(b) By Lemma~\ref{lem:thgrows} with $\sT-\sH=\alpha\ind{\cdot=h}$,
\begin{multline*}
  2\row_{r+1}-2\row_r\\
  =\alpha\bigl(\ind{r=h-1}-\ind{r=n-1-h}\bigr)+\sG(|2r+3-n|)-\sG(|2r+1-n|).
\end{multline*}
If $\alpha=0$ this vanishes. Let $2h\ge n$. One checks $|2r+1-n|\le2h-n-1$ iff $n-h\le r\le h-1$; hence
$\sG(|2r+3-n|)-\sG(|2r+1-n|)=\alpha(\ind{n-h-1\le r\le h-2}-\ind{n-h\le r\le h-1})=\alpha(\ind{r=n-h-1}-\ind{r=h-1})$
(both intervals are empty if $2h=n$). So all rows are equal, and Lemma~\ref{lem:rowvalue} gives their
value.

(c) By Lemma~\ref{lem:levels} it suffices to check \eqref{eq:cellTH} and \eqref{eq:cellTG}. For
$0\le s<u\le n-1$ with $u-s$ odd put $x=n-s$, $y=n-u$: then $1\le y<x\le n$, $x-y$ is odd, and
$\sT(s)+\sH(u)=\vartheta(x)+\vartheta(y)-\alpha\ind{y=K+1}$. For $s\le n-1$ with $n-s$ odd, $x=n-s$ is
odd and $\sT(s)+\sG(s)=\vartheta(x)+\alpha\ind{x\ge2K+3}$, because $s\le2h-n-1$ iff $x\ge2n-2h+1=2K+3$.
\end{proof}

In Lean: \lean{dipQ\_succ}, \lean{dip\_run}, \lean{dip\_row\_succ}, \lean{isIntervalUnif\_dip}
(\lean{DipoleUnif.lean}). It remains to choose $\alpha$ so that \eqref{eq:cell1} and \eqref{eq:cell2}
hold. For this we need a closed form of $\vartheta$, which involves the following special functions.

%==================================================================================================
\section{Special functions and the facts F1--F6}\label{sec:special}
%==================================================================================================

\begin{definition}\label{def:special}
For integers $x\ge1$, $K\ge0$ and $k\ge1$ let
\[
  \cB(x)=\sum_{j\ge0}\frac{(-1)^j}{x+j},\qquad \sigma(x)=4\cB(x)-\frac2x,\qquad
  a(K)=4K^2\cB(K+1)+1-2K,
\]
\[
  \Omega(k)=\sum_{y=k+1}^{2k}(-1)^y\frac{k}{y(y-1)}\in\mathbb{Q},\qquad w(K)=(-1)^K\,\Omega(K+1).
\]
\end{definition}

The series defining $\cB$ converges by the alternating series test. $\cB$ is Nielsen's beta function
at integer arguments: $\cB(1)=\ln2$ and $\cB(m+1)=(-1)^m(\ln2-\sum_{i=1}^m(-1)^{i-1}/i)$, which gives
the formula for $a(K)$ in the introduction; for instance $a(1)=3-4\ln2\approx0.2274$,
$a(2)=16\ln2-11\approx0.0904$, $a(3)=25-36\ln2\approx0.0467$. Also $\Omega(2)=-\frac16$,
$\Omega(3)=\frac15$, $\Omega(4)=-\frac{19}{210}$. These closed forms are not used below. In Lean,
$\cB$, $\sigma$, $a$, $\Omega$, $w$ are \lean{betaF} (defined as the limit of the partial sums),
\lean{sigmaF}, \lean{aF}, \lean{omegaF}, \lean{wF} (\lean{UnifFacts.lean}).

\begin{lemma}[alternating recursions]\label{lem:alt}
Let $x_0\in\N$ and $f,c:\{x_0,x_0+1,\dots\}\to\R$ with $f(x)+f(x+1)=c(x)$ for all $x\ge x_0$, $c$
nonincreasing and $f(x)\to0$. Then $0\le f(x)\le c(x)$ for all $x\ge x_0$. If moreover $x\mapsto c(x)-c(x+1)$
is nonincreasing, then
\[
  \tfrac12c(x)\le f(x)\le\tfrac12c(x)+\tfrac12\bigl(c(x)-c(x+1)\bigr).
\]
\end{lemma}

\begin{proof}
$f(x+2)-f(x)=c(x+1)-c(x)\le0$, so $N\mapsto f(x+2N)$ is nonincreasing with limit $0$; hence $f(x)\ge0$, and
$f(x)=c(x)-f(x+1)\le c(x)$. For the second part apply the first to $g(x)=f(x)-f(x+1)$, which satisfies
$g(x)+g(x+1)=c(x)-c(x+1)$ and $g\to0$: then $0\le g(x)\le c(x)-c(x+1)$, and $2f(x)=c(x)+g(x)$.
\end{proof}

\begin{lemma}\label{lem:B}
For integers $x\ge1$: $\cB(x)+\cB(x+1)=\frac1x$, $0\le\cB(x)\le\frac1x$, and $\cB(x)\to0$.
\end{lemma}

\begin{proof}
The partial sums satisfy $\sum_{j<N+1}\frac{(-1)^j}{x+j}=\frac1x-\sum_{j<N}\frac{(-1)^j}{x+1+j}$; let
$N\to\infty$. The partial sums of even length are $\le\cB(x)$ and those of odd length are $\ge\cB(x)$,
which gives $0\le\cB(x)\le1/x$.
\end{proof}

\begin{lemma}[Lemma B]\label{lem:lemmaB}
Let $u(x)=\frac1{2x}+\frac1{4x^2}-\frac1{8x^4}+\frac1{4x^6}$ and
$\varepsilon(x)=\frac{17x^4+34x^3+29x^2+12x+2}{8x^6(x+1)^6}$. For every integer $x\ge1$,
\[
  u(x)-\varepsilon(x)\le\cB(x)\le u(x) .
\]
\end{lemma}

\begin{proof}
A direct computation gives $u(x)+u(x+1)-\frac1x=\varepsilon(x)$ and
$\varepsilon(x)-\varepsilon(x+1)=\frac{17x^4+68x^3+100x^2+64x+16}{x^6(x+1)(x+2)^6}>0$. So
$f=u-\cB$ satisfies $f(x)+f(x+1)=\varepsilon(x)$ (Lemma~\ref{lem:B}) with $\varepsilon$ nonincreasing, and
$f\to0$. Lemma~\ref{lem:alt} gives $0\le f\le\varepsilon$.
\end{proof}

\begin{lemma}[properties of $\sigma$]\label{lem:sigma}
Let $x,M\ge1$ and $0\le K\le M$ be integers.
\begin{enumerate}
\item[(i)] $\sigma(x)+\sigma(x+1)=\frac2{x(x+1)}$.
\item[(ii)] $\frac1{x(x+1)}\le\sigma(x)\le\frac{x+4}{x(x+1)(x+2)}$; in particular $\sigma(x)>0$,
$\sigma(4)\le\frac1{15}$ and $\sigma(5)\le\frac3{70}$.
\item[(iii)] $\sigma(x+1)<\sigma(x)$.
\item[(iv)] $a(K)=\frac2{K+1}-1+K^2\sigma(K+1)$, and $a(M)=1-M^2\sigma(M)$.
\item[(v)] $M^2\sigma(M+1)\le1-\frac1{M+1}$.
\item[(vi)] $(M^2-K^2)\sigma(x)\le1-\frac1{M+1}$ for $x\ge M+1$.
\item[(vii)] $(M^2-K^2)\sigma(M+1)+a(M)\le1-K^2\sigma(M+1)$.
\item[(viii)] $1-K^2\sigma(x)\ge\frac2{K+1}-a(K)$ for $x\ge K+1$.
\end{enumerate}
\end{lemma}

\begin{proof}
(i) follows from $\cB(x)+\cB(x+1)=1/x$. (ii) Apply the second part of Lemma~\ref{lem:alt} to
$f=\sigma$, $c(x)=\frac2{x(x+1)}$, for which $c(x)-c(x+1)=\frac4{x(x+1)(x+2)}$ is nonincreasing, and
$\sigma\to0$. (iii) By (ii) at $x+1$ and since $x(x+5)<(x+2)(x+3)$,
$\sigma(x+1)\le\frac{x+5}{(x+1)(x+2)(x+3)}<\frac1{x(x+1)}$; by (i),
$\sigma(x)=\frac2{x(x+1)}-\sigma(x+1)>\frac1{x(x+1)}>\sigma(x+1)$. (iv) The first identity is the
definition rewritten with $\sigma(K+1)=4\cB(K+1)-\frac2{K+1}$; the second follows from the first (at $K=M$)
and (i) at $x=M$. (v) By (ii),
$M^2\sigma(M+1)\le\frac{M^2(M+5)}{(M+1)(M+2)(M+3)}\le\frac M{M+1}$, since $M(M+5)\le(M+2)(M+3)$.
(vi) By (iii), $0<\sigma(x)\le\sigma(M+1)$, and $0\le M^2-K^2\le M^2$; use (v). (vii) By (iv),
the left side is $2M^2\sigma(M+1)-K^2\sigma(M+1)+\frac2{M+1}-1$; use (v). (viii) By (iii) and (iv),
$1-K^2\sigma(x)\ge1-K^2\sigma(K+1)=\frac2{K+1}-a(K)$.
\end{proof}

\begin{lemma}\label{lem:omega}
For $k\ge1$, $\Omega(k)=(-1)^k\bigl(2k\cB(k+1)-1\bigr)+2k\cB(2k)-\frac12$. Consequently, for $K\ge0$,
\[
  w(K)=\begin{cases}
  \tfrac32-2(K+1)\bigl(\cB(K+2)+\cB(2K+2)\bigr), & K\text{ odd},\\[2pt]
  \tfrac12-2(K+1)\cB(K+2)+2(K+1)\cB(2K+2), & K\text{ even}.
  \end{cases}
\]
\end{lemma}

\begin{proof}
Write $\frac k{y(y-1)}=k(\frac1{y-1}-\frac1y)$, so
$\Omega(k)=-k\bigl(\sum_{z=k}^{2k-1}\frac{(-1)^z}z+\sum_{y=k+1}^{2k}\frac{(-1)^y}y\bigr)$. By
Lemma~\ref{lem:B}, $\sum_{z=a}^{a+N-1}\frac{(-1)^z}{z}=(-1)^a\cB(a)-(-1)^{a+N}\cB(a+N)$. Hence
$\Omega(k)=-k\bigl((-1)^k(\cB(k)-\cB(k+1))-\cB(2k)+\cB(2k+1)\bigr)$, and $\cB(k)=\frac1k-\cB(k+1)$,
$\cB(2k+1)=\frac1{2k}-\cB(2k)$ give the formula. For $w(K)$ take $k=K+1$ and multiply by $(-1)^K$.
\end{proof}

\begin{proposition}[facts F1--F6]\label{prop:facts}
\begin{enumerate}
\item[F1] $a(K)>0$ for all $K\ge0$.
\item[F2] $a(K+1)<a(K)$ for all $K\ge0$.
\item[F3] $a(K)+a(K+1)\le\frac1{2(K+1)}$ for all $K\ge2$.
\item[F4] $w(K)>0$ and $(K+1)\bigl(a(K)+a(K+1)\bigr)\le4w(K)$ for all $K\ge2$.
\item[F5] $a(1)-\frac2{21}-a(3)>0$, $\;a(1)-\frac2{21}+a(2)\le\frac27$, and $a(1)<\frac5{21}$.
\item[F6] $a(1)<\frac14$ and $a(3)<\frac1{20}$.
\end{enumerate}
\end{proposition}

\begin{proof}
Put $\aup(K)=4K^2u(K+1)+1-2K$ and $\alo(K)=4K^2\bigl(u(K+1)-\varepsilon(K+1)\bigr)+1-2K$. By
Lemma~\ref{lem:lemmaB} at $x=K+1$, $\alo(K)\le a(K)\le\aup(K)$ for all $K\ge0$.

F1: $a(K)\ge\alo(K)$, which is a quotient of polynomials with positive coefficients by \eqref{eq:A1}.

F2: $a(K)-a(K+1)\ge\alo(K)-\aup(K+1)>0$ by \eqref{eq:A2}.

F3: for $K=k+2$, $\frac1{2(K+1)}-a(K)-a(K+1)\ge\frac1{2(K+1)}-\aup(K)-\aup(K+1)\ge0$ by \eqref{eq:A3}.

F4: let $K=k+2\ge2$. If $K$ is odd, Lemma~\ref{lem:omega} and the upper bounds of Lemma~\ref{lem:lemmaB}
at $K+2$ and $2K+2$ give $w(K)\ge w^-_{\mathrm{odd}}(K):=\frac32-2(K+1)\bigl(u(K+2)+u(2K+2)\bigr)$,
which is positive by \eqref{eq:A4}; and by \eqref{eq:A5},
\[
  (K+1)\bigl(a(K)+a(K+1)\bigr)\le(K+1)\bigl(\aup(K)+\aup(K+1)\bigr)\le4w^-_{\mathrm{odd}}(K)\le4w(K).
\]
If $K$ is even, the upper bound at $K+2$ and the lower bound at $2K+2$ give
$w(K)\ge w^-_{\mathrm{even}}(K):=\frac12-2(K+1)u(K+2)+2(K+1)\bigl(u(2K+2)-\varepsilon(2K+2)\bigr)$, which is
positive by \eqref{eq:A6}, and \eqref{eq:A7} gives the second inequality in the same way.

F5, F6: Lemma~\ref{lem:lemmaB} at $x=2,3,4$ gives exact rational enclosures. With
$u(2)=\frac{79}{256}$, $u(3)=\frac{1127}{5832}$, $u(4)=\frac{2297}{16384}$ and the values of $\varepsilon$,
\begin{gather*}
  a(1)\in\Bigl[\tfrac{15}{64}-\tfrac{343}{46656},\ \tfrac{15}{64}\Bigr],\qquad
  a(2)\in\Bigl[\tfrac{67}{729}-\tfrac{1297}{746496},\ \tfrac{67}{729}\Bigr],\\
  a(3)\in\Bigl[\tfrac{193}{4096}-\tfrac{31689}{64000000},\ \tfrac{193}{4096}\Bigr],
\end{gather*}
that is $0.22702\le a(1)\le0.23438$, $0.09016\le a(2)\le0.09191$ and $0.04662\le a(3)\le0.04712$. Hence
$a(1)-\frac2{21}-a(3)\ge0.0846$, $a(1)-\frac2{21}+a(2)\le0.2311<\frac27$, $a(1)\le\frac{15}{64}<\frac5{21}$,
$a(1)<\frac14$ and $a(3)\le\frac{193}{4096}<\frac1{20}$.
\end{proof}

The polynomial identities \eqref{eq:A1}--\eqref{eq:A7} are listed in Appendix~\ref{app:poly}; they are
identities of rational functions, which can be checked by clearing denominators and expanding. In
Lean they are proved by \lean{field\_simp; ring} (\lean{poly\_F1} to \lean{poly\_F4d}), and F1--F6 are
\lean{factF1} to \lean{factF6} (\lean{UnifFacts.lean}); the Lean proof of F5, F6 uses the slightly weaker
decimal enclosures $0.227<a(1)<0.2344$, $0.0901<a(2)<0.092$, $0.0466<a(3)<0.0472$.

\begin{corollary}\label{cor:facts}
\textup{(a)} $a(L)\le a(K)$ for $L\ge K\ge0$. \textup{(b)} $a(K)<\frac1{2(K+1)}$ for $K\ge1$.
\end{corollary}

\begin{proof}
(a) is F2. (b) For $K=1$ this is F6; for $K\ge2$ it follows from F3 and $a(K+1)>0$ (F1).
\end{proof}

%==================================================================================================
\section{Closed form of \texorpdfstring{$\vartheta$}{theta}}\label{sec:closed}
%==================================================================================================

Throughout this section we are in the setting of Definition~\ref{def:dipole} with any $\alpha\in\R$.
Let $\Gamma=M^2-K^2$ and, for integers $x,y\ge1$,
\[
  \Phi(x)=\begin{cases}1-K^2\sigma(x), & x\le M,\\ \Gamma\,\sigma(x), & x\ge M+1,\end{cases}\qquad
  b(y)=\begin{cases}\dfrac{2(K+1)\alpha}{y(y-1)}, & K+2\le y\le2K+2,\\[4pt] 0, & \text{otherwise},\end{cases}
\]
and $\psi(x)=\vartheta(x)-\Phi(x)$ (the ``alternating part'').

\begin{lemma}\label{lem:closed}
\textup{(a)} $\vartheta(x)=0$ for $0\le x\le K$.
\textup{(b)} $\vartheta(K+1)=\frac2{K+1}$ and $\psi(K+1)=a(K)$.
\textup{(c)} For $x\ge K+1$, $\psi(x+1)=-\psi(x)+b(x+1)-\ind{x=M}\,a(M)$.
\end{lemma}

\begin{proof}
(a) $\cQ(x)=0$ for $x\le K$, since $x<K+1<M+1$ and $x<K+2$; so $\vartheta(x)=-\vartheta(x-1)=0$ for
$2\le x\le K$ by induction. (b)
$\cQ(K+1)=q_K(K+1)=2(K+1)K-2K^2=2K$, so $\vartheta(K+1)=\frac{2K}{(K+1)K}-\vartheta(K)=\frac2{K+1}$; then
$\psi(K+1)=\frac2{K+1}-1+K^2\sigma(K+1)=a(K)$ by Lemma~\ref{lem:sigma}(iv). (c) For $x\ge K+1$,
$\vartheta(x+1)+\vartheta(x)=\cQ(x+1)/(x(x+1))$ with
\[
  \cQ(x+1)=\begin{cases}2x(x+1)-2K^2+b(x+1)x(x+1), & x+1\le M,\\ 2M^2-2K^2+b(x+1)x(x+1), & x+1\ge M+1.\end{cases}
\]
If $x+1\le M$, Lemma~\ref{lem:sigma}(i) gives $\Phi(x+1)+\Phi(x)=2-\frac{2K^2}{x(x+1)}$, so
$\psi(x+1)+\psi(x)=b(x+1)$. If $x\ge M+1$, $\Phi(x+1)+\Phi(x)=\frac{2\Gamma}{x(x+1)}$ and again
$\psi(x+1)+\psi(x)=b(x+1)$. If $x=M$,
$\Phi(M+1)+\Phi(M)=\Gamma\sigma(M+1)+1-K^2\sigma(M)=\frac{2\Gamma}{M(M+1)}+1-M^2\sigma(M)
=\frac{2\Gamma}{M(M+1)}+a(M)$ by Lemma~\ref{lem:sigma}(i),(iv), so $\psi(M+1)+\psi(M)=b(M+1)-a(M)$.
\end{proof}

Unwinding the recursion, for $x\ge K+1$,
\begin{equation}\label{eq:psiclosed}
  \psi(x)=(-1)^x\Bigl((-1)^{K+1}a(K)+\sum_{y=K+2}^{x}(-1)^yb(y)+\ind{x\ge M+1}(-1)^Ma(M)\Bigr).
\end{equation}

\begin{remark}[the parity mode and the dipole]\label{rem:parity}
For $x\ge\max(M+1,2K+2)$, \eqref{eq:psiclosed} and $\sum_{y=K+2}^{2K+2}(-1)^yb(y)=2\alpha\,\Omega(K+1)$ give
$\vartheta(x)=\Gamma\sigma(x)+(-1)^x\bigl(A_{\mathrm{ren}}+2\alpha\,\Omega(K+1)\bigr)$ with
$A_{\mathrm{ren}}=(-1)^{K+1}a(K)+(-1)^Ma(M)$, which is nonzero because $a(K)>a(M)>0$. Without a dipole
($\alpha=0$) the smooth part $\Gamma\sigma(x)\approx\Gamma/x^2$ eventually becomes smaller than the parity
amplitude $|A_{\mathrm{ren}}|$, and cells can become negative when $n$ is large compared with $K$: for
example, for $(n,\ell,h)=(12,10,10)$ the matrix $Y^0=\cR^{10}_{12}-\cR^{9}_{12}$ has minimal entry
$-677/4620$, attained at the central cells $(5,6)$ and $(6,5)$. (That $Y^0=\cR^h_n-\cR^{\ell-1}_n$ in
general is not needed; it was checked for $n\le20$, Appendix~\ref{app:checks}.) The dipole amplitude
\begin{equation}\label{eq:alphastar}
  \alpha^*=\alpha^*(K,M)=-\frac{A_{\mathrm{ren}}}{2\,\Omega(K+1)}=\frac{a(K)-(-1)^{K+M}a(M)}{2\,w(K)}
\end{equation}
cancels this parity mode exactly. (In Lean: \lean{alphaStar}.)
\end{remark}

In Lean, Lemma~\ref{lem:closed} is \lean{dipTau\_low}, \lean{dipTau\_base}, \lean{dipS\_base} and
\lean{dipS\_succ} (\lean{UnifAllN.lean}).

%==================================================================================================
\section{The case analysis}\label{sec:cases}
%==================================================================================================

Throughout, $1\le\ell\le h\le n-2$, $K=n-1-h$, $M=n-\ell$, and $\vartheta$, $\Phi$, $\psi$, $b$ are as in
Sections~\ref{sec:dipole}--\ref{sec:closed}. Every pair $(n,[\ell,h])$ falls into exactly one of the cases
(C) $n\le2K+1$; (D) $n\ge2K+2$ and $K=1$; (E) $n\ge2K+2$ and $K\ge2$. We repeatedly use that
$(-1)^{x+c}=-(-1)^{y+c}$ when $x-y$ is odd, and Corollary~\ref{cor:facts}: for $M\ge K+1$,
$0<a(M)\le a(K+1)<a(K)$.

\subsection{Case (C): no dipole}

\begin{proposition}\label{prop:caseC}
If $n\le2K+1$, then $Y^0$ satisfies \eqref{eq:cell1} and \eqref{eq:cell2}; hence it is an interval
uniformisation of $[\ell,h]$.
\end{proposition}

\begin{proof}
If $K=1$ then $n\le3$, which forces $(n,\ell,h)=(3,1,1)$ and $M=2$. Here
$\vartheta(0),\dots,\vartheta(3)=0,0,1,0$ ($\cQ(2)=2$, $\cQ(3)=q_1(3)-q_2(3)=10-4=6$), the pairs in
\eqref{eq:cell1} are $(x,y)=(2,1),(3,2)$ with sums $1,1$, and $\vartheta(1)=\vartheta(3)=0$.

Let $K\ge2$ and $\alpha=0$, so $b\equiv0$. By Lemma~\ref{lem:closed} and induction, for $x\ge K+1$,
\[
  \vartheta(x)=\Phi(x)+(-1)^{x+K+1}a(K)+\ind{x\ge M+1}(-1)^{x+M}a(M).
\]
\emph{Bounds for single $x\le n$.} If $x\le K$, $\vartheta(x)=0$. If $K+1\le x\le M$, by
Lemma~\ref{lem:sigma}(viii) and Corollary~\ref{cor:facts}(b),
$\vartheta(x)\ge\frac2{K+1}-2a(K)>\frac2{K+1}-\frac1{K+1}>0$. If $M+1\le x\le n\le2K+1$, then
$\Gamma\ge(K+1)^2-K^2=2K+1$ and, by Lemma~\ref{lem:sigma}(ii),(iii),
$\sigma(x)\ge\sigma(2K+1)\ge\frac1{(2K+1)(2K+2)}$, so
$\vartheta(x)\ge\Gamma\sigma(x)-a(K)-a(K+1)\ge\frac1{2(K+1)}-a(K)-a(K+1)\ge0$ by F3. For the upper
bound, $\Phi(x)\le1$ (for $x\le M$ since $\sigma>0$; for $x>M$ by Lemma~\ref{lem:sigma}(vi)), so
$\vartheta(x)\le1+a(K)+a(K+1)\le1+\frac1{2(K+1)}<2$ by F3.

\emph{Pairs.} Let $K+1\le y<x\le n$ with $x-y$ odd. The $a(K)$ terms cancel. If $x\le M$:
$\vartheta(x)+\vartheta(y)=2-K^2(\sigma(x)+\sigma(y))\le2$. If $y\le M<x$: using
$\Gamma\sigma(x)\le\Gamma\sigma(M+1)$ and Lemma~\ref{lem:sigma}(vii),
$\vartheta(x)+\vartheta(y)\le\Gamma\sigma(M+1)+a(M)+1-K^2\sigma(y)\le2-K^2\sigma(M+1)-K^2\sigma(y)\le2$.
If $M<y$: the $a(M)$ terms cancel too, and $\vartheta(x)+\vartheta(y)=\Gamma(\sigma(x)+\sigma(y))\le2-\frac2{M+1}$
by Lemma~\ref{lem:sigma}(vi).

\emph{Conclusion.} In \eqref{eq:cell1} (with $\alpha=0$) both terms are $\ge0$; if $y\le K$ then
$\vartheta(y)=0$ and $\vartheta(x)\le2$, otherwise the pair bound applies. \eqref{eq:cell2} is
$0\le\vartheta(x)\le2$.
\end{proof}

In Lean: \lean{dipS\_C}, \lean{caseC\_bounds}, \lean{caseC\_cells}, \lean{caseC\_K1\_cells}.

\subsection{Case (E): one dipole of amplitude \texorpdfstring{$\alpha^*$}{alpha*}}

\begin{lemma}\label{lem:tail}
Let $\alpha\ge0$ and $\tail(x)=\sum_{j\ge0}(-1)^jb(x+1+j)$ (a finite sum). Then
$\tail(x)=b(x+1)-\tail(x+1)$, $\tail(x)=0$ for $x\ge2K+2$, $\tail(K+1)=2\alpha\,w(K)$, and
$0\le\tail(x)\le b(x+1)$ for $x\ge K+1$.
\end{lemma}

\begin{proof}
The first two statements are immediate. With $y=K+2+j$,
$\tail(K+1)=\sum_{j=0}^{K}(-1)^j\frac{2(K+1)\alpha}{(K+2+j)(K+1+j)}$, while
$w(K)=(-1)^K\sum_{y=K+2}^{2K+2}(-1)^y\frac{K+1}{y(y-1)}=\sum_{j=0}^{K}(-1)^j\frac{K+1}{(K+2+j)(K+1+j)}$.
For $x\ge K+1$ the sequence $b(x+1),b(x+2),\dots$ is nonnegative and nonincreasing, so its alternating
sum lies in $[0,b(x+1)]$.
\end{proof}

\begin{lemma}\label{lem:alphastar}
If $K\ge2$ and $M\ge K+1$, then $0<\alpha^*\le\frac2{K+1}$.
\end{lemma}

\begin{proof}
By F4, $w(K)>0$. The numerator of \eqref{eq:alphastar} is $\ge a(K)-a(M)>0$ and
$\le a(K)+a(M)\le a(K)+a(K+1)\le\frac{4w(K)}{K+1}$ by F4.
\end{proof}

\begin{proposition}\label{prop:caseE}
If $K\ge2$ and $n\ge2K+2$, then $Y^{\alpha^*}$ satisfies \eqref{eq:cell1} and \eqref{eq:cell2}; hence it
is an interval uniformisation of $[\ell,h]$.
\end{proposition}

\begin{proof}
Write $\alpha=\alpha^*$. By Proposition~\ref{prop:dipole}(b) the condition $2h\ge n$ is needed for the
rows; it is equivalent to $n\ge2K+2$. We show the cell inequalities for all $x,y\ge1$, not only
$x\le n$.

\emph{Closed form.} For $x\ge K+1$,
\begin{equation}\label{eq:psiE}
  \psi(x)=\tail(x)-\ind{x\le M}(-1)^{x+M}a(M).
\end{equation}
For $x=K+1$: by Lemma~\ref{lem:tail} and \eqref{eq:alphastar}, the right side is
$2\alpha w(K)+(-1)^{K+M}a(M)=a(K)=\psi(K+1)$. If \eqref{eq:psiE} holds at $x\ge K+1$, then by
Lemma~\ref{lem:closed}(c) and Lemma~\ref{lem:tail},
$\psi(x+1)=\tail(x+1)+\ind{x\le M}(-1)^{x+M}a(M)-\ind{x=M}a(M)$, which is \eqref{eq:psiE} at $x+1$ (for
$x<M$ directly; for $x=M$ the two $a(M)$ terms cancel; for $x>M$ both vanish).

\emph{Auxiliary bounds.} By Lemma~\ref{lem:alphastar}, $(K+1)\alpha\le2$, so for $z\ge K+1$,
$\tail(z)\le b(z+1)\le\frac4{z(z+1)}\le K^2\sigma(z)$ (Lemma~\ref{lem:sigma}(ii) and $K^2\ge4$) and
$\tail(z)\le\frac4{12}=\frac13$ (as $z\ge3$). Also $a(M)\le a(1)<\frac14$ and $a(M)\le a(K+1)<a(K)<\frac1{2(K+1)}$.

\emph{Lower bound $\vartheta\ge0$.} For $x\le K$, $\vartheta(x)=0$. For $K+1\le x\le M$, by
Lemma~\ref{lem:sigma}(viii), $\vartheta(x)\ge\frac2{K+1}-a(K)+0-a(M)>\frac2{K+1}-2a(K)>0$. For $x\ge M+1$,
$\vartheta(x)=\Gamma\sigma(x)+\tail(x)\ge0$.

\emph{Upper bounds.} $\Phi(x)\le1$ for $x\ge K+1$, so $\vartheta(x)\le1+\frac13+\frac14<2$ for all $x$. For
$x\ge2K+3$, $\tail(x)=0$ and $\alpha\le\frac2{K+1}\le\frac23$, so $\vartheta(x)+\alpha\le1+\frac14+\frac23<2$.

\emph{Pairs.} Let $K+1\le y<x$ with $x-y$ odd; we show $\vartheta(x)+\vartheta(y)\le2$. If $x\le M$, the
$a(M)$ terms cancel and $\vartheta(x)+\vartheta(y)=2-K^2\sigma(x)+\tail(x)-K^2\sigma(y)+\tail(y)\le2$. If
$y\le M<x$, then $\tail(x)\le\frac4{(M+1)(M+2)}\le K^2\sigma(M+1)$ and $\tail(y)\le K^2\sigma(y)$, so by
Lemma~\ref{lem:sigma}(vii),
\begin{align*}
  \vartheta(x)+\vartheta(y)&\le\bigl(\Gamma\sigma(M+1)+a(M)\bigr)+\tail(x)+1-K^2\sigma(y)+\tail(y)\\
  &\le\bigl(1-K^2\sigma(M+1)\bigr)+K^2\sigma(M+1)+1=2 .
\end{align*}
If $M<y$, then $\vartheta(x)+\vartheta(y)\le2\bigl(1-\frac1{M+1}\bigr)+\frac4{x(x+1)}+\frac4{y(y+1)}\le2$,
since $x,y\ge M+1\ge4$ gives $\frac4{z(z+1)}\le\frac1{M+1}$ for $z\in\{x,y\}$.

\emph{Conclusion.} \eqref{eq:cell1}: if $y=K+1$, then $\vartheta(y)-\alpha=\frac2{K+1}-\alpha\ge0$ gives the
lower bound, and $\vartheta(x)+\vartheta(K+1)-\alpha\le\vartheta(x)+\vartheta(K+1)\le2$ by the pair bound.
If $y\ne K+1$, the lower bound is $\vartheta\ge0$; the upper bound is the single bound if $y\le K$ and
the pair bound if $y\ge K+2$. \eqref{eq:cell2}: for $x\ge2K+3$ use $\vartheta(x)+\alpha\le2$, otherwise
$0\le\vartheta(x)<2$.
\end{proof}

In Lean: \lean{alphaStar\_bounds}, \lean{dipTail\_base}, \lean{dipTail\_bounds}, \lean{dipS\_E},
\lean{caseE\_cells\_gen}, \lean{caseE\_cells}.

\subsection{Case (D): \texorpdfstring{$K=1$}{K=1} and one dipole of amplitude \texorpdfstring{$2/7$}{2/7}}

\begin{proposition}\label{prop:caseD}
If $K=1$ (i.e.\ $h=n-2$) and $n\ge4$, then $Y^{2/7}$ satisfies \eqref{eq:cell1} and \eqref{eq:cell2};
hence it is an interval uniformisation of $[\ell,h]$.
\end{proposition}

\begin{proof}
Let $\alpha=\frac27$; then $2h=2n-4\ge n$. Here $2\le M\le n-1$, $\Gamma=M^2-1$, $b(3)=\frac4{21}$,
$b(4)=\frac2{21}$ and $b(y)=0$ for $y\ge5$. Put $e_P=a(1)-\frac2{21}$ and $A=e_P+(-1)^Ma(M)$.

\emph{Values.} $\vartheta(1)=0$ and $\vartheta(2)=1$ (Lemma~\ref{lem:closed}). By Lemma~\ref{lem:closed},
$\psi(2)=a(1)$, $\psi(3)=-a(1)+\frac4{21}-\ind{M=2}a(2)$, and for $x\ge4$
\[
  \psi(x)=(-1)^x\bigl(e_P+\ind{x\ge M+1}(-1)^Ma(M)\bigr)
\]
(check $x=4$ for $M=2$, $M=3$, $M\ge4$ separately, then induct with $b(x+1)=0$). Since $a(1)=\sigma(2)$,
$a(2)=4\sigma(3)-\frac13$ and $\sigma(2)+\sigma(3)=\frac13$ (Lemma~\ref{lem:sigma}(i),(iv)), we get
$\vartheta(3)=\frac67$ if $M\ge3$ and $\vartheta(3)=\frac4{21}$ if $M=2$. For $x\ge4$:
$\vartheta(x)=1-\sigma(x)+(-1)^xe_P$ if $x\le M$, and $\vartheta(x)=\Gamma\sigma(x)+(-1)^xA$ if $x\ge M+1$.

\emph{Constants.} By F5 and F1, $0<a(1)-\frac2{21}-a(3)<e_P<\frac5{21}-\frac2{21}=\frac17$. If $M$ is even,
$A=e_P+a(M)\in(0,e_P+a(2)]\subseteq(0,\frac27]$ by F2 and F5; if $M$ is odd, then $M\ge3$ and
$A=e_P-a(M)\ge e_P-a(3)>0$, $A<e_P<\frac27$. By Lemma~\ref{lem:sigma}(ii),(iii), $\sigma(x)\le\frac1{15}$ for
$x\ge4$ and $\sigma(x)\le\frac3{70}$ for $x\ge5$. For $x\ge M+1$: $0\le\Gamma\sigma(x)\le1-\frac1{M+1}$ and
$\Gamma\sigma(x)+a(M)\le\Gamma\sigma(M+1)+a(M)\le1-\sigma(M+1)<1$ (Lemma~\ref{lem:sigma}(vi),(vii)).

\emph{Single $x\ge4$.} Even $x$: $0<\vartheta(x)\le1+e_P$ (for $x\le M$, $\vartheta(x)\ge\frac{14}{15}$; for
$x>M$, $\vartheta(x)=\Gamma\sigma(x)+A>0$ and $\Gamma\sigma(x)+A\le\Gamma\sigma(x)+a(M)+e_P<1+e_P$). Odd $x$:
$-\frac27\le\vartheta(x)\le1$ (for $x\le M$, $x\ge5$ and $\vartheta(x)=1-\sigma(x)-e_P\in[1-\frac3{70}-\frac17,1]$;
for $x>M$, $\vartheta(x)=\Gamma\sigma(x)-A\in[-A,1-A]$).

\emph{Pairs $4\le y<x$, $x-y$ odd.} If $x\le M$: $\vartheta(x)+\vartheta(y)=2-\sigma(x)-\sigma(y)\in[0,2]$. If
$y\le M<x$: $\vartheta(x)+\vartheta(y)=\Gamma\sigma(x)+1-\sigma(y)-(-1)^{y+M}a(M)$, which is
$\le\Gamma\sigma(M+1)+a(M)+1-\sigma(y)\le2$ and $\ge1-\frac1{15}-a(2)>0$. If $M<y$:
$\vartheta(x)+\vartheta(y)=\Gamma(\sigma(x)+\sigma(y))\in[0,2]$.

\emph{Conclusion.} \eqref{eq:cell1} with $y\ge4$ is the pair bound ($y\ne K+1=2$). For $y=1$ ($x$ even):
$\vartheta(x)\in\{1\}\cup[0,1+e_P]$. For $y=2$ ($x$ odd, $\vartheta(2)-\alpha=\frac57$):
$\vartheta(3)+\frac57\in\{\frac{19}{21},\frac{11}7\}$, and $\vartheta(x)+\frac57\in[\frac37,\frac{12}7]$ for
$x\ge5$. For $y=3$ ($x\ge4$ even): $0\le\vartheta(x)+\vartheta(3)\le1+e_P+\frac67<2$. \eqref{eq:cell2}:
$\vartheta(1)=0$, $\vartheta(3)\in[0,2]$, and for odd $x\ge5$ (where $x\ge2K+3$),
$\vartheta(x)+\frac27\in[0,\frac97]$.
\end{proof}

In Lean: \lean{dipB\_D3}, \lean{dipB\_D4}, \lean{dipS\_D}, \lean{tauD\_three}, \lean{caseD\_cells}.

\subsection{Proof of \texorpdfstring{Theorem~\ref{thm:unif}}{Theorem 3}}\label{sec:proofunif}
Let $n\ge1$ and $0\le\ell\le h\le n-1$. Case (A), $h=n-1$: take $J_n$ if $\ell=0$ and
$J_n-\cR^{\ell-1}_n$ if $\ell\ge1$ (Lemma~\ref{lem:compl}). Case (B), $\ell=0$ and $h\le n-2$: take the
renewal square $\cR^h_n$ (Proposition~\ref{prop:renewal}). Otherwise $1\le\ell\le h\le n-2$, and
Propositions~\ref{prop:caseC}, \ref{prop:caseD} and~\ref{prop:caseE} cover the three remaining cases
(C), (D), (E). \qed

In Lean this is \lean{intervalUnifAll\_all} (\lean{UnifAllN.lean}).

%==================================================================================================
\section{Formal verification}\label{sec:lean}
%==================================================================================================

\subsection{What is formalised}
Theorem~\ref{thm:diag} is formalised, with its complete proof, in Lean~4 \cite{deMouraUllrich2021}
using Mathlib \cite{mathlib2020}. The development is the library \lean{STUProof} in the directory
\lean{lean/} of the repository \cite{STUrepo} (\url{https://github.com/anshM123/IQOQI-Conjecture}). The
final theorem is \lean{STUProof.stu\_exists\_unconditional} in \lean{KirwanFreeFinal.lean}; its statement
is displayed in Remark~\ref{rem:leanstatement}. All lemmas, propositions and theorems of
Sections~\ref{sec:reduction}--\ref{sec:cases} are formalised (Lemma~\ref{lem:hall} in the inequality form
stated after it), including Lemma~\ref{lem:horn} (Horn's theorem), the analytic estimates of
Section~\ref{sec:special} for all $K$ (not by finite computation), and the polynomial identities of
Appendix~\ref{app:poly}. The function $\cB$ is defined as the limit of its partial sums, whose existence
comes from Mathlib's alternating series test.

The reduction of Theorem~\ref{thm:main} to Theorem~\ref{thm:diag} is \lean{stu\_exists\_hermitian}
(\lean{GeneralH.lean}). Not formalised are the physical discussion of Section~\ref{sec:intro}, the closed
forms involving $\ln2$ mentioned after Definition~\ref{def:special} (they are not used), the remarks of
Section~\ref{sec:remarks} (except for the ingredients of the Kirwan route, see Section~\ref{sec:kirwan}),
and the side remarks marked as checked by computer in Appendix~\ref{app:checks}.

\subsection{Versions and checking}
The toolchain is Lean \lean{v4.33.1} (file \lean{lean-toolchain}), and Mathlib is pinned to commit
\lean{0df444a360eaa60ab8c11dca51a86af692955474} (file \lean{lakefile.toml}). The script
\lean{check.sh} elaborates all files sequentially (to bound memory), then runs \lean{Axioms.lean}, which
prints the axioms used by $67$ main declarations, and finally searches the sources for the keywords
\lean{sorry}, \lean{admit}, \lean{native\_decide} and for \lean{axiom} declarations. A complete run
(log \lean{logs/lean\_check\_final\_unconditional.log} in \cite{STUrepo}) finished with exit code $0$
after $1573$ seconds (about $26$ minutes) and a peak memory use of roughly $7$\,GB. For every one of the
printed declarations, including \lean{stu\_exists\_unconditional}, the output has the form
\begin{alltt}\small
'STUProof.stu_exists_unconditional' depends on axioms:
    [propext, Classical.choice, Quot.sound]
\end{alltt}
that is, only Lean's three standard axioms (propositional extensionality, choice, and quotient
soundness) are used, and the keyword search reports no occurrence. In particular no \lean{sorry}, no
added axiom and no \lean{native\_decide} (which would trust compiled code) is used.

\subsection{Trusted base}
Besides the Lean kernel, a reader who wants to trust the formal result must check that the statement
means what it should. It uses Mathlib's real exponential, finite sums, matrices, the Kronecker product
and the unitary group, and four definitions of the project, displayed in
Remark~\ref{rem:leanstatement}: \lean{gibbs}, \lean{gibbsState}, \lean{partialTraceA} and
\lean{partialTraceB}. No other definition of the project enters the statement.

\subsection{Correspondence}
Table~\ref{tab:lean} lists the main results of this paper and their Lean counterparts (all in the
namespace \lean{STUProof}).

\begin{table}[ht]
\centering\small
\begin{tabular}{p{0.32\textwidth}p{0.38\textwidth}p{0.2\textwidth}}
\toprule
Paper & Lean declaration(s) & File\\
\midrule
Theorem~\ref{thm:diag} & \lean{stu\_exists\_unconditional} & \lean{KirwanFreeFinal}\\
Proposition~\ref{prop:reduction} & \lean{stu\_exists\_of\_intervalUnif} & \lean{KirwanFreeU}\\
Theorem~\ref{thm:unif} & \lean{intervalUnifAll\_all} & \lean{UnifAllN}\\
Eq.~\eqref{eq:localunitary} & \lean{partialTraceB\_conj\_kronecker}, \lean{partialTraceA\_conj\_kronecker} & \lean{STU}\\
Lemma~\ref{lem:dom} & \lean{SMaj}, \lean{SMaj.add}, \lean{SMaj.smul}, \lean{SMaj.majOn}, \lean{smaj\_self}, \lean{smaj\_zeroOne} & \lean{KirwanFreeCore}\\
Lemma~\ref{lem:horn} & \lean{schurHorn} & \lean{SchurHorn}\\
Proposition~\ref{prop:realisation} & \lean{ReachTri.orthogonal}, \lean{tri\_reachable}, \lean{lc\_reachable} & \lean{KirwanFreeCore}\\
Lemma~\ref{lem:convex} & \lean{ReachTri.convex} & \lean{KirwanFreeCore}\\
Lemma~\ref{lem:T}, Corollary~\ref{cor:thermal} & \lean{lemmaT}, \lean{thermal\_clip\_weights} & \lean{KirwanFreeCore}\\
Lemmas~\ref{lem:telescope}, \ref{lem:decomp} & \lean{clipZ\_eq}, \lean{one\_sub\_clipZ}, \lean{wk\_cell\_identity} & \lean{KirwanFree}\\
Lemma~\ref{lem:crossrun} & \lean{crossInd\_run}, \lean{card\_filter\_Icc\_hook}, \lean{crossInd\_sum} & \lean{KirwanFreeHook}\\
Lemma~\ref{lem:hall} (inequality form) & \lean{leftM\_nonneg}, \lean{Ltot\_pos}, \lean{sum\_aM} & \lean{KirwanFree}\\
Definition~\ref{def:unif}, Lemma~\ref{lem:unifdom} & \lean{IsIntervalUnif}, \lean{IntervalUnifAll}, \lean{IsIntervalUnif.smaj} & \lean{IntervalUnif}\\
Proposition~\ref{prop:wk} & \lean{wkU\_reach}, \lean{wkU\_smaj}, \lean{wkU\_rows} & \lean{KirwanFreeU}\\
Proposition~\ref{prop:renewal}, Lemma~\ref{lem:compl} & \lean{isIntervalUnif\_renU}, \lean{isIntervalUnif\_boxJ}, \lean{IsIntervalUnif.compl} & \lean{Renewal}\\
Lemmas~\ref{lem:levels}--\ref{lem:rowvalue} & \lean{thg\_cells}, \lean{thg\_run\_rec}, \lean{thg\_row\_succ}, \lean{row\_const\_value} & \lean{DipoleUnif}\\
Lemma~\ref{lem:Qdiff}, Proposition~\ref{prop:dipole} & \lean{dipQ\_succ}, \lean{dip\_run}, \lean{dip\_row\_succ}, \lean{isIntervalUnif\_dip} & \lean{DipoleUnif}\\
Lemmas~\ref{lem:alt}--\ref{lem:omega} & \lean{alt\_rec\_bounds}, \lean{alt\_rec\_bounds2}, \lean{betaF\_add}, \lean{betaF\_bounds}, \lean{lemmaB}, \lean{sigmaF\_bounds}, \lean{sigmaF\_succ\_lt}, \lean{aF\_eq\_sigma}, \lean{omegaF\_eq}, \lean{wF\_eq} & \lean{UnifFacts}\\
Lemma~\ref{lem:sigma}(v)--(viii) & \lean{msig\_le}, \lean{gamma\_sig\_le}, \lean{gamma\_sig\_add\_a}, \lean{phi\_low\_ge} & \lean{UnifAllN}\\
Proposition~\ref{prop:facts} & \lean{factF1}, \dots, \lean{factF6} & \lean{UnifFacts}\\
Lemma~\ref{lem:closed} & \lean{dipTau\_low}, \lean{dipTau\_base}, \lean{dipS\_base}, \lean{dipS\_succ} & \lean{UnifAllN}\\
Propositions~\ref{prop:caseC}, \ref{prop:caseE}, \ref{prop:caseD} & \lean{caseC\_cells}, \lean{caseC\_K1\_cells}, \lean{caseE\_cells}, \lean{caseD\_cells} & \lean{UnifAllN}\\
\bottomrule
\end{tabular}
\caption{Paper results and Lean declarations (files in \lean{lean/STUProof/}, extension \lean{.lean}).}
\label{tab:lean}
\end{table}

\subsection{Other content of the repository}
The repository also contains formalised results that are not needed for Theorem~\ref{thm:diag}: the
Kirwan route of Section~\ref{sec:kirwan}, with Kirwan's theorem as an explicit hypothesis; unconditional
proofs for $d\le20$ and $d\le40$ from kernel-checked certificates of ``hook decompositions'' (a stronger
combinatorial statement that implies $\mathrm{UNIF}(n)$); and an alternative explicit ``cross''
construction for interval uniformisations. These are superseded by Theorem~\ref{thm:diag}.

%==================================================================================================
\section{Remarks}\label{sec:remarks}
%==================================================================================================

\subsection{The strong majorization form is false for \texorpdfstring{$d\ge4$}{d >= 4}}
For $d=3$, every $q\prec p$ is symmetrically reachable from $p\otimes p$
\cite[Thm.~1]{BakhshinezhadEtAl2019}. This fails in every higher dimension, so the thermal structure of
the target (Lemma~\ref{lem:T}) cannot be dispensed with.

\begin{proposition}\label{prop:counterexample}
Let $d\ge4$, $\delta=\frac1{5d}$, $\eta=2\delta^2$, and
\begin{gather*}
  p=(1-\delta-(d-2)\eta,\ \delta,\ \eta,\dots,\eta),\qquad s=1-p_0=\delta+(d-2)\eta,\\
  q=\Bigl(p_0,\ \frac s{d-1},\dots,\frac s{d-1}\Bigr).
\end{gather*}
Then $p$ is strictly positive and nonincreasing (so it is the Gibbs vector of $H=\diag(-\ln p_i)$ at
$\beta=1$), $q\prec p$, but no state $\rho$ in the unitary orbit of $\diag(p\otimes p)$ has both marginals
with spectrum $q$. For $d=4$: $p=(0.94,0.05,0.005,0.005)$ and $q=(0.94,0.02,0.02,0.02)$.
\end{proposition}

\begin{proof}
Since $\delta\le\frac1{20}$ and $2(d-2)\delta<\frac25$, we have $s\le\frac75\delta$ and $p_0>\frac12$, so
$p_0>\delta>\eta>0$; $q$ is nonincreasing, and $q\prec p$ because $q$ is the image of $p$ under the doubly
stochastic matrix $1\oplus\frac1{d-1}J_{d-1}$. The eigenvalues of $\diag(p\otimes p)$, in nonincreasing
order $\lambda_1\ge\lambda_2\ge\cdots$, start with
\[
  \lambda_1=p_0^2>\lambda_2=\lambda_3=p_0\delta>\lambda_4=\dots=\lambda_{2d-1}=p_0\eta ,
\]
and all others are products $p_ip_j$ with $i,j\ge1$, hence $\le\delta^2<p_0\eta$ (as $p_0>\frac12$). Note
$2\lambda_1+\lambda_2+\dots+\lambda_{2d-1}=2p_0(p_0+\delta+(d-2)\eta)=2p_0$, and the remaining eigenvalues
sum to $s^2$.

Suppose $\rho$ is in the orbit and both marginals have spectrum $q$. Applying a local unitary $V\otimes W$
(which preserves the orbit and the marginal spectra) we may assume that $|0\rangle$ is an eigenvector of
$\rho_A$ and of $\rho_B$ for the eigenvalue $p_0$. Let $P_2=|00\rangle\langle00|$ and let $P_1$ be the
projection onto $O=\mathrm{span}\{|0j\rangle,|j0\rangle:j\ge1\}$, of rank $2(d-1)$. With
$X=|0\rangle\langle0|\otimes1+1\otimes|0\rangle\langle0|=2P_2+P_1$,
\begin{align*}
  2p_0&=\langle0|\rho_A|0\rangle+\langle0|\rho_B|0\rangle=\Tr[X\rho]=\Tr[P_2\rho]+\Tr[(P_2+P_1)\rho]\\
  &\le\lambda_1+(\lambda_1+\dots+\lambda_{2d-1})=2p_0 ,
\end{align*}
by Ky Fan's maximum principle \cite{Bhatia1997}. So both Ky Fan bounds are equalities. If a
rank-$r$ projection $P$ satisfies $\Tr[P\rho]=\lambda_1+\dots+\lambda_r$ and $\lambda_r>\lambda_{r+1}$, then
$P$ is the spectral projection of $\rho$ for its $r$ largest eigenvalues (writing
$\Tr[P\rho]=\sum_i\lambda_i\langle v_i|P|v_i\rangle$ with weights in $[0,1]$ summing to $r$, equality forces
$\langle v_i|P|v_i\rangle=1$ for $i\le r$). Because of the gaps $\lambda_1>\lambda_2$ and
$\lambda_{2d-1}>\lambda_{2d}$, $\rho$ is block diagonal,
$\rho=\lambda_1P_2\oplus\rho_O\oplus\rho_R$ with respect to
$\C|00\rangle\oplus O\oplus R$, $R=\mathrm{span}\{|ij\rangle:i,j\ge1\}$, where $\rho_O$ has eigenvalues
$\lambda_2,\dots,\lambda_{2d-1}$ and $\Tr\rho_R=s^2$.

Now compute $\rho_A$. The partial trace over $B$ of $|0j\rangle\langle j'0|$ vanishes for $j,j'\ge1$, and
$\Tr_B\rho_R$ is supported on $\mathrm{span}\{|j\rangle:j\ge1\}$. Hence
$\rho_A=\langle0|\rho_A|0\rangle\,|0\rangle\langle0|\oplus(Y+Z)$, where $Y$ is the compression of $\rho_O$
to $\mathrm{span}\{|j0\rangle:j\ge1\}\cong\C^{d-1}$ and $Z=\Tr_B\rho_R\ge0$, $\Tr Z=s^2$. Since
$\langle0|\rho_A|0\rangle=p_0$ and $\rho_A$ has spectrum $q$, all $d-1\ge3$ eigenvalues of $Y+Z$ equal
$s/(d-1)$. But by Weyl's inequality and Cauchy interlacing \cite{Bhatia1997},
\begin{align*}
  \lambda_3(Y+Z)&\le\lambda_3(Y)+\lambda_1(Z)\le\lambda_3(\rho_O)+\Tr Z=p_0\eta+s^2\\
  &<2\delta^2+2\delta^2=\frac{4\delta}{5d}<\frac{\delta}{d-1}\le\frac s{d-1},
\end{align*}
using $s^2\le\frac{49}{25}\delta^2$. This contradiction proves the claim.
\end{proof}

For $d=4$ the final inequality reads $0.0083<0.02$. The inequalities in this proof were also checked in
exact arithmetic for $4\le d\le400$ (Appendix~\ref{app:checks}).

Conversely, the set of symmetric marginal spectra is not contained in $\{q:q\prec p\}$ either. For
$d=3$ and $p=(0.4,0.35,0.25)$, mixing the three cells $(1,1)$, $(0,2)$, $(2,0)$ (distinct rows and
columns) to the common diagonal value $0.1075$, which is majorized by $(p_1^2,p_0p_2,p_2p_0)$, gives
(Lemma~\ref{lem:horn} and the argument of Proposition~\ref{prop:realisation}) a state whose two marginals
are both $\diag(0.4075,0.335,0.2575)$; its largest entry exceeds $p_0$.

\subsection{The Kirwan route}\label{sec:kirwan}
We sketch an alternative proof of Theorem~\ref{thm:diag}, which uses Kirwan's convexity theorem and is
\emph{not needed} for the results of this paper. For $k\in[d]$ let
$v_k=(\frac{P_k}{k+1},\dots,\frac{P_k}{k+1},p_{k+1},\dots,p_{d-1})$ with $P_k=p_0+\dots+p_k$.
\begin{enumerate}
\item Every $v_k$ is a symmetric marginal spectrum of $p\otimes p$, via an explicit real orthogonal
matrix: the tail levels $j>k$ are relabelled by cyclic shifts of the top block $T=\{0,\dots,k\}$ (chosen
by greedy bin packing, so that each shift class carries weight at most $\frac1{k+1}$), and on each cyclic
diagonal $\{(a,a+s)\}\subseteq T\times T$ a Schur--Horn rotation realises the image of the diagonal
under one doubly stochastic circulant matrix. Lean: \lean{vk\_reachable}, \lean{exists\_binPacking}
(\lean{Construction.lean}).
\item $\gamma_{\beta'}(E)\in\conv\{v_0,\dots,v_{d-1}\}$ (Lemma~5 of \cite{BakhshinezhadEtAl2019}; Lean:
\lean{lemma5}, \lean{Lemma5.lean}).
\item By Kirwan's theorem \cite{Kirwan1984}, applied to $U(d)\times U(d)$ acting on the unitary orbit of
$\diag(p\otimes p)$ with moment map $\rho\mapsto(\Tr_B\rho,\Tr_A\rho)$, the set of pairs of sorted marginal
spectra is a convex polytope \cite{Klyachko2004,ChristandlMitchison2006}; its intersection with the
diagonal $\{(q,q)\}$ is convex and contains all $v_k$, hence $\gamma_{\beta'}(E)$.
\end{enumerate}
In Lean, step~3 is the hypothesis \lean{KirwanSymmetricConvexity d}, and \lean{stu\_exists}
(\lean{STU.lean}) proves Theorem~\ref{thm:diag} from it. Kirwan's theorem itself is not formalised.

\subsection{Further remarks}
\emph{More general targets.} The argument proves more than Theorem~\ref{thm:diag}: every
$q\in\conv\{w_0,\dots,w_{d-1}\}$ belongs to $\Reach(p)$ and is therefore a symmetric marginal of
$\diag(p\otimes p)$. By Lemma~\ref{lem:T} this includes $g(p)/\sum_ig(p_i)$ for every concave nondecreasing
$g:[0,\infty)\to[0,\infty)$ with $\sum_ig(p_i)>0$. This generalisation is not stated separately in Lean,
but it follows from the formalised statements \lean{lemmaT}, \lean{wkU\_reach}, \lean{ReachTri.convex}
and \lean{ReachTri.orthogonal} by the argument of Proposition~\ref{prop:reduction}.

\emph{Explicitness.} The configuration $F$ is explicit: its entries are rational functions of $p$ and
of the numbers $a(K)$, $w(K)$ (which involve $\ln 2$). The orthogonal matrix is block diagonal over the
$2d-1$ diagonals of the array, and on each diagonal it is a product of Givens rotations and permutations
obtained constructively from Lemma~\ref{lem:horn} \cite{ChanLi1983}.

\emph{Numerical confirmation.} Independently of the proof, the whole chain was implemented from the
formulas of this paper: for $300$ random Hermitian $H$ with $d\le6$ (including degenerate spectra) the
resulting unitaries have both marginals equal to $\tau_{\beta'}(H)$ up to $4\cdot10^{-15}$
(Appendix~\ref{app:checks}).

\subsection{Outlook}
The set $\mathcal S(p)$ of symmetric marginal spectra of $\diag(p\otimes p)$ remains to be understood: we
know $\conv\{w_k\}\cup\conv\{v_k\}\subseteq\mathcal S(p)$ (the second inclusion via Kirwan), that
$\mathcal S(p)\ne\{q:q\prec p\}$ for $d\ge4$, and that $\mathcal S(p)\not\subseteq\{q:q\prec p\}$. Natural
next questions are the analogous problem for more than two copies, quantitative (robust) versions, and
whether the fixed diagonal structure used here also settles other quantum marginal problems with
diagonal marginals.


%==================================================================================================
\appendix
\section{Polynomial identities for F1--F4}\label{app:poly}
%==================================================================================================

With $u,\varepsilon$ from Lemma~\ref{lem:lemmaB} and $\aup,\alo,w^-_{\mathrm{odd}},w^-_{\mathrm{even}}$
from the proof of Proposition~\ref{prop:facts}, the following are identities of rational functions (in
\eqref{eq:A3}, \eqref{eq:A5} and \eqref{eq:A7}, $K=k+2$). All numerator coefficients are positive.
\begin{align}
  \alo(K)&=\frac{K^5+13K^4+70K^3+194K^2+256K+128}{2(K+1)(K+2)^6},\label{eq:A1}\\
  \alo(K)-\aup(K+1)&=\frac{2K^4+23K^3+101K^2+170K+98}{2(K+1)(K+2)^6},\label{eq:A2}\\
  \frac1{2(K+1)}-\aup(K)-\aup(K+1)&=\frac{P_3(k)}{2(k+3)^6(k+4)^7},\label{eq:A3}\\
  w^-_{\mathrm{odd}}(K)&=\frac{P_{4a}(K)}{128(K+1)^5(K+2)^6},\label{eq:A4}\\
  \frac{4w^-_{\mathrm{odd}}(K)}{K+1}-\aup(K)-\aup(K+1)&=\frac{P_{4b}(k)}{32(k+3)^6(k+4)^6},\label{eq:A5}\\
  w^-_{\mathrm{even}}(K)&=\frac{P_{4c}(K)}{4(K+2)^6(2K+3)^6},\label{eq:A6}\\
  \frac{4w^-_{\mathrm{even}}(K)}{K+1}-\aup(K)-\aup(K+1)&=\frac{P_{4d}(k)}{2(k+3)^6(k+4)^6(2k+7)^6},\label{eq:A7}
\end{align}
where
\begin{align*}
P_3(k)&=k^{12}+41k^{11}+764k^{10}+8545k^9+63789k^8+334153k^7+1256064k^6+3400638k^5\\
&\quad+6544384k^4+8654345k^3+7354236k^2+3503150k+660280,\\
P_{4a}(K)&=48K^{10}+768K^9+5442K^8+22428K^7+59385K^6+105468K^5+127292K^4\\
&\quad+103200K^3+53840K^2+16320K+2176,\\
P_{4b}(k)&=16k^{10}+544k^9+8242k^8+73212k^7+421633k^6+1641880k^5+4368048k^4\\
&\quad+7817152k^3+8974576k^2+5940736k+1710240,\\
P_{4c}(K)&=160K^{11}+3200K^{10}+29012K^9+157296K^8+566317K^7+1420999K^6+2535030K^5\\
&\quad+3215114K^4+2841125K^3+1666285K^2+583870K+92626,\\
P_{4d}(k)&=192k^{16}+10816k^{15}+285192k^{14}+4672168k^{13}+53227666k^{12}+447135728k^{11}\\
&\quad+2864967224k^{10}+14282437922k^9+55985185735k^8+173131057544k^7\\
&\quad+420984335806k^6+796458028404k^5+1149337523504k^4+1223024950168k^3\\
&\quad+905103455591k^2+416231111464k+89613810730.
\end{align*}
These are the identities \lean{poly\_F1}, \lean{poly\_F2}, \lean{poly\_F3}, \lean{poly\_F4a},
\lean{poly\_F4b}, \lean{poly\_F4c}, \lean{poly\_F4d} of \lean{UnifFacts.lean}, where they are proved by
\lean{field\_simp; ring}. We also verified them independently with exact computer algebra
(Appendix~\ref{app:checks}).

%==================================================================================================
\section{Independent computer checks}\label{app:checks}
%==================================================================================================

The proof does not depend on any computation beyond the formal verification. To guard against
transcription errors between the Lean development and this paper, every formula stated here was
re-implemented from the paper and checked; the scripts and their logs are available with the source
files of this paper (directory \lean{checks/}).
\begin{itemize}
\item \lean{check\_identities.py} (exact computer algebra): the ingredients of Lemma~\ref{lem:lemmaB},
the identities of Lemmas~\ref{lem:sigma}, \ref{lem:omega} and \ref{lem:Qdiff}, the renewal identities of
Proposition~\ref{prop:renewal}, the Hall identity of Lemma~\ref{lem:hall}, and
\eqref{eq:A1}--\eqref{eq:A7} with positivity of all coefficients.
\item \lean{check\_facts\_numeric.py} and \lean{check\_small\_values.py}: F1--F6 and
Corollary~\ref{cor:facts} at $60$ digits for $K\le3000$ (F4 for $K\le1200$), the exact enclosures of
$a(1),a(2),a(3)$, the closed forms $a(1)=3-4\ln2$, etc., $0<\alpha^*\le\frac2{K+1}$ for
$2\le K\le150$, $K+1\le M\le K+150$, and agreement of the amplitudes of Table~\ref{tab:cases} with an
independent implementation in \cite{STUrepo} (all intervals with $n\le40$).
\item \lean{check\_constructions.py}: the renewal squares agree with the block description ($n\le24$);
exact verification of Definition~\ref{def:unif} for all intervals of all boxes $n\le22$ with the
constructions of Table~\ref{tab:cases} ($\alpha^*$ rounded to a rational with denominator
$\le10^{30}$); agreement of $\sT$ with an independent implementation of the run recursion ($n\le30$); the
cell conditions \eqref{eq:cell1}, \eqref{eq:cell2} with the exact $\alpha^*$ at $50$ digits for all
$17296$ intervals with $1\le\ell\le h\le n-2$, $n\le48$; the closed forms of Section~\ref{sec:closed};
$Y^0=\cR^h_n-\cR^{\ell-1}_n$ for $n\le20$; and floating-point checks near the top ($K\le6$, $n\le400$)
and for random intervals with $n$ up to $2000$.
\item \lean{check\_thg\_generic.py}: Lemmas~\ref{lem:levels}--\ref{lem:rowvalue} with random rational
$\sT,\sH,\sG$.
\item \lean{check\_wk.py}: Lemmas~\ref{lem:telescope}--\ref{lem:hall} and Proposition~\ref{prop:wk} in
exact arithmetic for random rational spectra (with ties) with $d\le8$ and every $k$; Lemma~\ref{lem:T}
and Corollary~\ref{cor:thermal} for $300$ random thermal instances.
\item \lean{check\_end\_to\_end.py}: the full construction of Theorem~\ref{thm:main}, including the
reduction for non-diagonal $H$, with explicit Givens rotations, for $300$ random instances with $d\le6$.
\item \lean{check\_counterexample.py}: the inequalities in the proof of
Proposition~\ref{prop:counterexample} for $4\le d\le400$ (exact), a numerical test of the interlacing
step, and the $d=3$ example.
\end{itemize}

\bibliographystyle{amsplain}
\bibliography{refs}

\end{document}
